\exercisesection{线性代数的补充}

\begin{enumerate}
\def\labelenumi{\arabic{enumi})}
\tightlist
\item
  设
\end{enumerate}

\[
\mathbf {M} \rightarrow \mathbf {N} \xrightarrow {t} \mathbf {P} \rightarrow \mathbf {Q}
\]

\[
\begin{array}{c c c c c} \downarrow u & & \downarrow v & & \downarrow w \\ \downarrow & & \downarrow & & \downarrow s \end{array}
\]

\[
\mathrm{M} ^ {\prime} \rightarrow \mathrm{N} ^ {\prime} \xrightarrow {t ^ {\prime}} \mathrm{P} ^ {\prime} \rightarrow \mathrm{Q} ^ {\prime}
\]

A-模的一个交换图，其中各行都是正合的；假设 \(u\) 是满射且 \(s\) 是单射。证明：

\[
\operatorname{Ker} (w) = t (\operatorname{Ker} (v)) \quad \operatorname{Im} (v) = t ^ {\prime - 1} (\operatorname{Im} (w)).
\]

\begin{enumerate}
\def\labelenumi{\arabic{enumi})}
\setcounter{enumi}{1}
\tightlist
\item
  设 \(u: \mathbf{M} \to \mathbf{N}, v: \mathbf{N} \to \mathbf{P}\) 为 A-模的两个同态。证明存在正合列：
\end{enumerate}

\[
0 \rightarrow \operatorname{Ker} (u) \rightarrow \operatorname{Ker} (v \circ u) \rightarrow \operatorname{Ker} (v) \rightarrow \operatorname{Coker} (u) \rightarrow \operatorname{Coker} (v \circ u) \rightarrow \operatorname{Coker} v \rightarrow 0.
\]

\begin{enumerate}
\def\labelenumi{\arabic{enumi})}
\setcounter{enumi}{2}
\tightlist
\item
  在 A-模的图中：
\end{enumerate}

\[
\begin{array}{c} \mathrm{P} \xleftarrow {} \mathrm{M} \\ \Big \downarrow \quad \mathrm{R} ^ {\swarrow} \Big \downarrow \\ \mathrm{N} \xleftarrow {} \mathrm{Q} \end{array}
\]

图中提取出的两个三角形的交换性是否必然蕴含整个图的交换性？

\begin{enumerate}
\def\labelenumi{\arabic{enumi})}
\setcounter{enumi}{3}
\item
  设 \(k\) 为一个域， \(B = k[T]\) ，设 A 为子环 \(k[T^2, T^3] \subset k[T]\) 。证明 A-模 \(B\) 是无挠的，但不是平坦 A-模。
\item
  设 \(k\) 为一个域， \(\mathbf{C} = k[\mathbf{X}, \mathbf{Y}]\) ；考虑子环 \(\mathbf{A} = k[\mathbf{X}^4, \mathbf{Y}^4]\) , \(\mathbf{B} = k[\mathbf{X}^4, \mathbf{X}^3\mathbf{Y}, \mathbf{XY}^3, \mathbf{Y}^4]\) 与 \(\tilde{\mathbf{B}} = k[\mathbf{X}^4, \mathbf{X}^3\mathbf{Y}, \mathbf{X}^2\mathbf{Y}^2, \mathbf{XY}^3, \mathbf{Y}^4]\) ，它们都是 \(\mathbf{C}\) 的子环。包含关系 \(\mathbf{A} \subset \mathbf{B} \subset \tilde{\mathbf{B}}\) 使我们可以在 \(\mathbf{B}\) 与 \(\tilde{\mathbf{B}}\) 上定义 A-模结构。证明 \(\tilde{\mathbf{B}}\) 是平坦 A-模，而 A-模 \(\mathbf{B}\) 不是平坦的。
\item
  证明：一个 A-模 E 是平坦的，当且仅当它满足以下条件（参见 X，第 15 页，注记）：
\end{enumerate}

(ii'') 对于 A 中元素的每个有限族 \((c_{i})_{i\in I}\) ，方程的每个解 \(e = (e_i)_{i\in I}\) （方程 \(\sum_{i\in I}c_i e_i = 0\) 的解）都可以写成 \(b_1 z_1 + \cdots + b_n z_n\) ，其中 \(b_1, \ldots, b_n \in E\) ，并且对于 \(r = 1, \ldots, n\) , \(z_r = (z_{r,i})_{i \in I}\) 是方程 \(\sum_{i=1}^{n} z_{r,i} c_i = 0\) 的一个解。

\begin{enumerate}
\def\labelenumi{\arabic{enumi})}
\setcounter{enumi}{6}
\tightlist
\item
  设 E 是一个幺半群，S 是 E 的一个满足消去律且满足 III，第 121 页，习题 17 的假设的子半群，因而存在一个“分式幺半群” \(\overline{\mathbf{E}}\) 以及一个单同态 \(\varepsilon: \mathbf{E} \to \overline{\mathbf{E}}\) 。设 A 是一个交换环。由 \(\varepsilon\) 可导出幺半群代数的一个同态 \(u: \mathbf{A}^{(\mathbf{E})} \to \mathbf{A}^{(\overline{\mathbf{E}})}\) ，它使我们可以把 \(\mathbf{A}^{\overline{\mathbf{E}}}\) 视为左 \(\mathbf{A}^{(\mathbf{E})}\) -模。证明这个模在 \(\mathbf{A}^{(\mathbf{E})}\) 上是平坦的（将 \(\mathbf{A}^{\overline{\mathbf{E}}}\) 写成秩为 1 的自由 \(\mathbf{A}^{(\mathbf{E})}\) -模的滤过归纳极限）。
\end{enumerate}

* 8) 设 A 为区间 {[}0, 1{]} 上连续实值函数构成的环。

\begin{enumerate}
\def\labelenumi{\alph{enumi})}
\item
  证明在 0 处取值为零的函数构成的理想是平坦 A-模（证明它是理想 Af 的滤过归纳极限，其中 f 是仅在 0 处取值为零的函数）。
\item
  证明在 0 的某个邻域中取值为零的函数构成的理想 I 是投射 A-模（证明存在一个序列 \((f_n)_{n\geqslant 1}\) , \(f_n\in \mathbf{I}\) ，使得 \(\operatorname {Supp}(f_n)\subset \left[\frac{1}{n + 2},\frac{1}{n}\right]\) ，并且 \(\sum_{n}f_{n}(x) = 1\) 对 \(x\neq 0\) 成立；由此推出
\end{enumerate}

I 的每个元素 \(h\) 都可以写成 \(h = \sum_{n}^{t_{h}^{\prime}} (f_n h) e_n\) ，其中 \(e_n = \sum_{r \leqslant n+1}^{n} f_r\) ，并使用 II，第 46 页，命题 12。*

\begin{enumerate}
\def\labelenumi{\arabic{enumi})}
\setcounter{enumi}{8}
\tightlist
\item
  设 P 是一个 A-模。证明以下条件是等价的：
\end{enumerate}

α) 为使右 A-模的序列 \(M' \xrightarrow{u} M \xrightarrow{v} M''\) 正合，必要且充分条件是序列

\[
\mathbf {M} ^ {\prime} \otimes_ {\mathbf {A}} \mathbf {P} \xrightarrow {u \otimes 1} \mathbf {M} \otimes_ {\mathbf {A}} \mathbf {P} \xrightarrow {v \otimes 1} \mathbf {M} ^ {\prime \prime} \otimes_ {\mathbf {A}} \mathbf {P}
\]

是正合的。

\(\beta)\) P 是平坦的，并且对于每一个非零右 A-模 M，有 M \(\otimes_{A} P \neq (0)\) .

γ) P 是平坦的，且对所有右 A 模 M、N，典范同态

\[
\operatorname{Hom} _ {\mathbf {A}} (\mathbf {M}, \mathbf {N}) \rightarrow \operatorname{Hom} _ {\mathbf {Z}} (\mathbf {M} \otimes_ {\mathbf {A}} \mathbf {P}, \mathbf {N} \otimes_ {\mathbf {A}} \mathbf {P})
\]

是单射。

δ) P 是平坦的，并且对于 A 的每一个极大右理想 m，都有 mP ≠ P。

此时称 P 为忠实平坦 A-模。

\begin{enumerate}
\def\labelenumi{\arabic{enumi})}
\setcounter{enumi}{9}
\tightlist
\item
  \begin{enumerate}
  \def\labelenumii{\alph{enumii})}
  \tightlist
  \item
    证明一个忠实平坦模是忠实的（II，第28页）。
  \end{enumerate}
\end{enumerate}

\begin{enumerate}
\def\labelenumi{\alph{enumi})}
\setcounter{enumi}{1}
\item
  Z-模 Q 是忠实且平坦的，但不是忠实平坦的。
\item
  设 \((p_n)\) 为严格递增的素数序列，并设 A 为乘积环 \(\prod_{n} \mathbf{Z}/p_n \mathbf{Z}\) 。
\end{enumerate}

证明 \(Z/p_{n}Z\) 的直和 E 是忠实的投射 A-模，但不是忠实平坦的（注意 E 是 A 的一个理想，满足 \(E^{2}=E\) ）。

\begin{enumerate}
\def\labelenumi{\arabic{enumi})}
\setcounter{enumi}{10}
\item
  设 \(0 \to \mathbf{P}' \to \mathbf{P} \to \mathbf{P}'' \to 0\) 为 A-模的正合列。假设 \(\mathbf{P}'\) 与 \(\mathbf{P}''\) 是平坦的，且其中一个是忠实平坦的。证明 \(\mathbf{P}\) 是忠实平坦的。
\item
  设 \(\rho: A \to B\) 为一个环同态，使得 \(B\) 成为 \(A\) 上的忠实平坦右模；设 \(M\) 为一个 \(A\) 模。
\end{enumerate}

\begin{enumerate}
\def\labelenumi{\alph{enumi})}
\item
  证明：M 是有限生成的（相应地，有限表示的），当且仅当 B-模 \(\mathbf{B} \otimes_{\mathbf{A}} \mathbf{M}\) 亦然。
\item
  我们假设 \(\rho(A)\) 包含在 B 的中心中。证明：M 是平坦的（相应地，忠实平坦的，相应地，有限生成投射的），当且仅当 B-模 \(B \otimes_{A} M\) 是如此。
\end{enumerate}

\begin{enumerate}
\def\labelenumi{\arabic{enumi})}
\setcounter{enumi}{12}
\tightlist
\item
  证明：整环 A 上每个有限生成的平坦模 P 都是投射的（考虑一个表示 \(0 \rightarrow R \xrightarrow{L} L \xrightarrow{P} 0\) ，其中 L 是有限生成的自由模；选取 R 的一个有限生成子模 \(R'\) ，使得 \(R/R'\) 是挠模，并应用第 14 页定理 1）。
\end{enumerate}

¶ 14) a) 设 \(0 \to \mathbf{R} \to \mathbf{L} \to \mathbf{E} \to 0\) 为 A-模的一个正合列，其中 L 是自由 A-模；设 \((e_{a})\) 为 L 的一组基。证明以下条件是等价的：

β) 对每个 \(x \in \mathbb{R}\) ，若 \(\mathfrak{a}_{x}\) 是 \(x\) 在基 \((e_{\alpha})\) 上的分量生成的右理想，则有 \(x \in \mathfrak{a}_{x} \mathbb{R}\) 。

\(\gamma)\) 对每一个 \(x \in \mathbb{R}\) ，存在同态 \(u_x: \mathbb{L} \to \mathbb{R}\) 使得 \(u_x(x) = x\) 。

δ) 对 R 的元素的每个有限序列 \((x_{i})_{1\leqslant i\leqslant n}\) ，存在同态 \(u:L\to R\) 使得 \(u(x_{i})=x_{i}\) 对 \(1\leqslant i\leqslant n\) 成立（利用第 14 页定理 1。）

\begin{enumerate}
\def\labelenumi{\alph{enumi})}
\setcounter{enumi}{1}
\item
  设 r 为 A 的根，并设 \(0 \rightarrow R \rightarrow L \rightarrow E \rightarrow 0\) 为 A-模的一个正合列，其中 L 是自由的。假设 E 是平坦的，且 R 包含于 rL 中。证明此时 R = 0（用 a) 的记号，注意 \(a_{x}\) 是有限生成的理想，并且我们有 \(a_{x} = a_{x} r\) ）。
\item
  设 E 是一个平坦的有限生成 A-模；假设存在 A 的一个双边理想 b 包含于 A 的根中，使得 E/bE 是自由 (A/b)-模。证明此时 E 是自由 A-模（注意到存在有限生成自由 A-模 L 使得 L/bL 同构于 E/bE，并应用 b））。特别地，局部环上的每个有限生成平坦模都是自由的。
\end{enumerate}

\begin{enumerate}
\def\labelenumi{\arabic{enumi})}
\setcounter{enumi}{14}
\tightlist
\item
  设 A 为一个环， \(a\) 为 A 的一个元素。证明以下性质是等价的： \(\alpha)\) \(a\in aAa\) 。
\end{enumerate}

β) Aa 是模 \(A_{s}\) 的一个直和项。

γ) \(A_{s}/Aa\) 是平坦的 A-模。

δ) 对于 A 的每个右理想 b，我们有 b ∩ Aa = ba。

\begin{enumerate}
\def\labelenumi{\arabic{enumi})}
\setcounter{enumi}{15}
\tightlist
\item
  设 A 是一个环。证明以下性质是等价的：
\end{enumerate}

α) 每个元素 \(a \in A\) 满足习题15的等价性质。

β) 每个有限生成的左理想都是 \(A_{s}\) 的直和项。

γ) 每个左 A-模都是平坦的。

δ) 每个右A-模都是平坦的。

我们于是称 A 为绝对平坦环（参见 VIII，§ 8，习题 15）。

\begin{enumerate}
\def\labelenumi{\arabic{enumi})}
\setcounter{enumi}{16}
\tightlist
\item
  设 A 是一个绝对平坦环（习题 16）。
\end{enumerate}

\begin{enumerate}
\def\labelenumi{\alph{enumi})}
\item
  设 P 为投射 A-模。证明 P 的每个有限生成子模都是 P 的直和项（化归到 P 为有限生成自由模的情形）。
\item
  证明每一个投射A-模都是循环子模的直和，且这些循环子模同构于A的循环理想。（利用Kaplansky定理（II，第183页，习题2）将情形化归为P具有可数生成族的情形，然后利用a）。）
\item
  证明 A 的每个由可数生成元族生成的理想都是投射的。
\item
  给出一个绝对平坦环 A 以及一个不是投射模的有限生成 A-模的例子。
\end{enumerate}

18) 设 M 是一个 A-模。证明以下条件是等价的：

α) M 是有限生成的（相应地，有限表示的）。

β) 对于相对于滤过预序集 I 的每个 A-模归纳系 \((\mathbf{N}_{i}, u_{ji})\) ，典范同态 \(\varinjlim\operatorname{Hom}_{\mathbf{A}}(\mathbf{M}, \mathbf{N}_{i}) \to \operatorname{Hom}_{\mathbf{A}}(\mathbf{M}, \varinjlim \mathbf{N}_{i})\) 是单射（相应地，双射）。

γ) 对于右 A-模的每一个族 \((\mathbf{P}_{j})_{j\in J}\) ，典范同态

\[
\left(\prod_ {j \in J} P _ {j}\right) \otimes_ {A} M \rightarrow \prod_ {j \in J} \left(P _ {j} \otimes_ {A} M\right)
\]

（II，第61页）是满射（相应地，双射）。

δ) 对于每个集合 I，典范同态 \(A_{d}^{1} \otimes_{A} M \to M^{1}\) 是满射（相应地，双射）。 19) 设 A 和 B 是两个环， \(\varphi : A \to B\) 为一个同态，M 是一个 B-模。我们赋予 \(B \otimes_{A} M\) 与 \(\operatorname{Hom}_{A}(B, M)\) 由 B 的 B-模结构诱导的左 B-模结构。称 M 相对于 A 是投射的，或 \(\varphi\) -投射的，如果典范满射 \(B \otimes_{A} M \to M\) 具有 B-线性截面；称 M 相对于 A 是内射的，或 \(\varphi\) -内射的，如果典范单射

\[
\mathbf {M} \rightarrow \operatorname{Hom} _ {\mathbf {A}} (\mathbf {B}, \mathbf {M})
\]

有一个 B-线性收缩。

\begin{enumerate}
\def\labelenumi{\alph{enumi})}
\item
  如果 M 作为 A-模是投射的（相应地，内射的）且是 \(\varphi\) -投射的（相应地， \(\varphi\) -内射的），则 M 是投射的（相应地，内射的）B-模。
\item
  对于每个 A-模 N，B-模 \(B \otimes_{A} N\) （相应地， \(\operatorname{Hom}_{A}(B, N)\) ）是 \(\varphi\) -投射的（相应地， \(\varphi\) -内射的）。
\item
  我们假设 \(\varphi(A)\) 包含在 B 的中心中。设 N 是一个 A-模，P 是一个 \(\varphi\) -投射的 B-模。证明 B-模 \(P \otimes_{A} N\) （相应地，右 B-模 \(\operatorname{Hom}_{A}(P, N)\) ）是 \(\varphi\) -投射的（相应地， \(\varphi\) -内射的）。
\end{enumerate}

\begin{enumerate}
\def\labelenumi{\arabic{enumi})}
\setcounter{enumi}{19}
\item
  设 A 是一个主理想环， \(a\) 为 A 的一个非零元素，B = A/Aa。证明 B-模 B \(_{s}\) 是内射的。
\item
  证明以下条件是等价的：
\end{enumerate}

α) A 是诺特环。

β) 每个作为内射 A-模的归纳极限的 A-模都是内射的。

γ) 每个作为内射 A-模的直和的 A-模都是内射的。

（为了证明 \(\gamma\) ）蕴含 \(\alpha\) ），设 \((\mathfrak{a}_n)_{n\geqslant 1}\) 为 A 的左理想组成的一个递增序列， \(\mathfrak{a} = \bigcup_{n}\mathfrak{a}_{n}\) ,

\(I_{n}\) 为 \(\mathfrak{a}/\mathfrak{a}_{n}\) 的一个内射包络。从 \(\mathfrak{a}\) 到 \(\oplus I_{n}\) 的自然同态是某个同态 \(f: A \to \oplus I_{n}\) 的限制；考虑 \(f(1)\) 。

\begin{enumerate}
\def\labelenumi{\arabic{enumi})}
\setcounter{enumi}{21}
\tightlist
\item
  设 A 为主理想环，K 为其分式域， \(\pi\) 为 A 的一个极值元素。证明 A-模 K/A 的 \(\pi\) -准素分量是 A-模 A/ \(\pi^n\) A 的一个内射包络，对每一个 \(n \geqslant 1\) 。 ¶ 23) 设 I 是一个内射 A-模，E 为环 End \(_A\) (I)，r 为 E 中使得 I 是 Ker( \(u\) ) 的内射包络的元素 \(u\) 组成的集合。
\end{enumerate}

\begin{enumerate}
\def\labelenumi{\alph{enumi})}
\item
  证明 r 是一个双边理想，并且环 E/r 是绝对平坦的（习题 16）。（设 u 是 E 的一个元素，N 是 I 的一个子模，使得 N ∩ Ker u = 0 且对此性质是极大的；证明存在 v ∈ E 使得对所有 x ∈ N 有 vu(x) = x，并由此推出 uvu - u ∈ r。）
\item
  证明 r 是 E 的根（包含关系 r(E) ⊂ r 由 a) 得出）；证明对于每个 u ∈ r，1 - u 是双射的）。
\item
  假设 I 是有限个不可分解内射模的直和。证明环 E/r 是半单的。
\end{enumerate}

设 \(i: \mathbf{A} \to \mathbf{I}\) 为 A-模 \(\mathbf{A}_{s}\) 的一个内射包络。我们设 \(\mathbf{E} = \operatorname{End}_{\mathbf{A}}(\mathbf{I})\) 与 \(\mathbf{B} = \operatorname{End}_{\mathbf{E}}(\mathbf{I})\) ，使得 \(\mathbf{B}\) 是 A-模 I 的双交换子（VIII，§ 3，定义 1）。我们把 \(\mathbf{A}\) 视为 \(\mathbf{B}\) 的一个子环。用 \(j: \mathbf{B} \to \mathbf{I}\) 与 \(p: \mathbf{E} \to \mathbf{I}\) 表示由 \(j(b) = b(i(1))\) 与 \(p(u) = u(i(1))\) （对 \(b \in \mathbf{B}\) , \(u \in \mathbf{E}\) ）定义的映射。

\begin{enumerate}
\def\labelenumi{\alph{enumi})}
\item
  证明 \(p\) 是满射且 \(j\) 是单射。
\item
  证明以下条件是等价的：
\end{enumerate}

α) 映射 p 是双射的。

β) 映射 j 是双射。

γ) A-模 B 是内射的。

若这些条件得到满足，则映射 \(j^{-1} \circ p\) 是从环 E 到 B 的相反环上的一个同构。

\begin{enumerate}
\def\labelenumi{\alph{enumi})}
\setcounter{enumi}{2}
\item
  设 \(\mathfrak{F}\) 为左理想 \(a\) （ \(A\) 中）组成的集合，满足 \(a \cap b \neq (0)\) 对每个非零左理想 \(b\) 都成立。用 \(s(A)\) 表示满足如下条件的元素 \(a\) （ \(A\) 中）组成的集合：存在 \(a \in \mathfrak{F}\) 使得 \(aa = 0\) 。证明 \(s(A)\) 是 \(A\) 的一个双边理想。
\item
  证明以下条件是等价的：
\end{enumerate}

α) \(s(\mathbf{A}) = 0.\)

β) 环 E 是半本原的。

γ) 环 B 是绝对平坦的。

若这些条件满足，则 \(b\) ) 的等价性质成立。

（设 \(s(I)\) 为元素 \(x\) （I 中）组成的集合，其中存在 \(\mathfrak{a} \in \mathfrak{F}\) 使得 \(\mathfrak{a} x = 0\) 。利用习题 23，证明 \(p^{-1}(s(I))\) 等于 E 的根 \(r\) ，由此得到蕴含关系 \(\beta) \Rightarrow \alpha\) 。然后证明 \(\alpha\) 蕴含 \(r = \text{Ker } (p) = 0\) ，它蕴含 \(\beta\) 与 \(\gamma\) （根据习题 23），也蕴含条件 \(\alpha\) （ \(b\) ) 的）。最后证明 \(\gamma\) 蕴含 \(\alpha\) 。）

\begin{enumerate}
\def\labelenumi{\alph{enumi})}
\setcounter{enumi}{4}
\item
  若 A 是左诺特的，则理想 s(A) 是幂零的（首先证明 s(A) 的每个元素 a 都是幂零的，考虑 \(a^{n}\) 的左零化子；然后利用 VIII，§ 9，习题 8 b））。¶ 25) 假设环 A 是左诺特的且不含非零幂零双边理想。a) 证明习题 24 中定义的环 B 是半单的（“Goldie 定理”；利用习题 23 c）以及习题 24 d）和 e））。若此外 (0) 是 A 中的素双边理想（VIII，§ 8，习题 19），证明 B 是单的。
\item
  将 A 等同于 B 的一个子环。证明以下性质：
\end{enumerate}

\begin{enumerate}
\def\labelenumi{(\roman{enumi})}
\item
  A 中每个左可消去的元素在 B 中都可逆（从而在 A 中可消去）。
\item
  每个元素 \(b\) （ \(\mathbf{B}\) 中）都可以写成 \(s^{-1}a\) ，其中 \(s \in \mathbf{A}\) , \(a \in \mathbf{A}\) ，且 \(s\) 可消去。
\end{enumerate}

有时称 B 为 A 的左全分式环。

（为证明 (ii)，归结到 B 是单的情形。然后证明 A 的理想 (0) 是素的。设 a 是 F 的一个理想，a 是 a 中使得右理想 aB 在理想 a'B（a' ∈ a）中极大的元素，x 是 B 中满足 axa = a 的元素。证明 B(1 - xa) ⊂ aB，由此

\[
(\mathbf {B} (1 - a x) \cap \mathfrak {a}) \cdot (\mathbf {B} (1 - x a) \cap \mathfrak {a}) = (0);
\]

推出 \(a\) 可消去。最后取 a 为所有 \(a \in A\) 使得 \(ab \in A\) 的元素组成的集合来得出结论。）c) 假设 \(A\) 的每个非零元素都是可消去的。证明 \(B\) 是一个域，它是 \(A\) 的左分式域（参见第一卷，第 155 页，习题 15）。

¶ 26) a) 证明以下条件是等价的：

α) 环 A 是右诺特的，且 A-模 \(A_{s}\) 是内射的。

\(\alpha^{0}\) A 是左诺特的且 A-模 \(A_{d}\) 是内射的。

β) 环 A 是自内射的（VIII，§ 2，习题 11）。

（证明 \(\alpha\) ) 蕴含每个右理想 r 满足 \(r = r^0\) 。然后证明 \(r(A)\) 是幂零的，并利用习题 23 推出 A 是右阿廷环。接着证明单左 A-模的对偶是单模。由此推出 A 是左诺特环，并且最终是自内射的。b) 证明环 A 是自内射的，当且仅当每个内射 A-模都是投射的且每个投射 A-模都是内射的。

¶ 27) 假设环 A 是左诺特的。

\begin{enumerate}
\def\labelenumi{\alph{enumi})}
\item
  设 I 是一个不可分解的内射 A-模。证明 I 的非零子模的零化子理想组成的集合具有一个最大元 \(\mathfrak{A}(I)\) ，它是一个素双边理想（VIII，§ 8，习题 19）。
\item
  证明：对于 A 的每一个素双边理想 p，存在一个不可分解的内射 A-模 I，使得 \(\mathfrak{A}(I) = \mathfrak{p}\) （考虑 A/p 的内射包络的一个不可分解直因子）。
\item
  假设 A 是交换的。证明每个不可分解的内射 A-模 I 都是 A/U(I) 的内射包络。由此推出映射 I ↦ U(I) 定义了从不可分解内射 A-模的同构类集合 S 到 A 的素理想集合上的一个双射。逆双射把理想 p 对应于 A/p 的内射包络。
\item
  证明：为使映射 I ↦ \(\mathfrak{A}(I)\) （从 \(\mathfrak{A}\) 到 A 的素双边理想集合上）是双射，只需 A 满足以下条件 (H)：
\end{enumerate}

\begin{enumerate}
\def\labelenumi{(\Alph{enumi})}
\setcounter{enumi}{7}
\tightlist
\item
  对于 A 的每一个左理想 a，存在 A/a 的元素 \(x_{1}, \ldots, x_{k}\) ，使得
\end{enumerate}

\[
\operatorname{Ann} (\mathbf {A} / \mathfrak {a}) = \bigcap_ {i} \operatorname{Ann} (x _ {i}).
\]

\begin{enumerate}
\def\labelenumi{\alph{enumi})}
\setcounter{enumi}{4}
\tightlist
\item
  证明在以下情形中条件 (H) 得到满足：
\end{enumerate}

α) A 的所有左理想都是双边理想。

β) A 是阿廷的。

γ) A 的中心 Z 是诺特环，且 A 是有限生成的 Z-模。

\begin{enumerate}
\def\labelenumi{\alph{enumi})}
\setcounter{enumi}{5}
\tightlist
\item
  设 k 为一个域， \(Q^{+}\) 为满足 \(\geqslant 0\) 的有理数组成的加法幺半群，A 为代数 \(k^{(\mathbf{Q}^{+})}\) ，a 为 A 的一个非零元素，I 为 A/Aa 的内射包络。证明 A-模 I 是不可分解的，但不包含与 A/p 同构的子模，其中 p 为素理想（注意 A 的理想集是全序的，且 A 仅包含两个素理想）。
\end{enumerate}

\begin{enumerate}
\def\labelenumi{\arabic{enumi})}
\setcounter{enumi}{27}
\tightlist
\item
  设 A 为诺特环，M 为 A-模，I 为 M 的内射包络，它是某个族 \((I_{\alpha})_{\alpha \in J}\) 的不可分解子模的直和。我们用 Ass(M) 表示素双边理想 \(\mathfrak{U}(I_{\alpha})\) （ \(\alpha \in J\) ）组成的集合（习题 27）；若 \(p \in \operatorname{Ass}(M)\) ，则称 p 与 M 相伴。
\end{enumerate}

\begin{enumerate}
\def\labelenumi{\alph{enumi})}
\item
  若 M ≠ 0，则 Ass (M) ≠ ∅；若 M 是有限生成的，则集合 Ass (M) 是有限的。
\item
  证明：素双边理想 p 与 M 相伴当且仅当 M 包含一个子模 N，使得对 N 的每个非零子模 N'，都有 Ann (N') = p。若 A 是交换的，则 p 与 M 相伴当且仅当 M 包含一个与 A/p 同构的子模。
\item
  设 M' 是 M 的一个子模。证明我们有
\end{enumerate}

\[
\operatorname{Ass} \left(\mathbf {M} ^ {\prime}\right) \subset \operatorname{Ass} (\mathbf {M}) \subset \operatorname{Ass} \left(\mathbf {M} ^ {\prime}\right) \cup \operatorname{Ass} \left(\mathbf {M} / \mathbf {M} ^ {\prime}\right).
\]

\begin{enumerate}
\def\labelenumi{\alph{enumi})}
\setcounter{enumi}{3}
\tightlist
\item
  我们设 \(J_{p} = \sum_{\mathfrak{M}(I_{\alpha}) = p} I_{\alpha}\) 与 \(Q(p) = \left( \sum_{q \neq p} J_{q} \right) \cap M\) 。证明我们有 \(\bigcap_{p \in \operatorname{Ass}(M)} Q(p) = 0\) ，并且对 Ass(M) 的每个异于 Ass(M) 的子集 E，有 \(\bigcap_{p \in E} Q(p) \neq 0\) 。对每个 \(p \in \operatorname{Ass}(M)\) ，我们有
\end{enumerate}

Ass \((\mathbf{M} / \mathbf{Q}(\mathfrak{p})) = \{\mathfrak{p}\}\) .

\begin{enumerate}
\def\labelenumi{\alph{enumi})}
\setcounter{enumi}{4}
\tightlist
\item
  设 a 为 A 的中心的一个元素。证明：以 a 为比率的位似变换在 M 中是单射的，其必要且充分条件是 a 不属于与 M 相伴的任何一个素理想。
\end{enumerate}

假设环 A 是交换的且是诺特的。设 a 为 A 的一个理想。

\begin{enumerate}
\def\labelenumi{\alph{enumi})}
\tightlist
\item
  设 H 为一个 A-模，使得 H 的每个元素都被 a 的某个幂零化。证明：要使 H 是内射的，只需对每个整数 \(n \geqslant 1\) 、每个被 \(a^{n}\) 零化的有限生成 A-模 M，以及 M 的每个子模 \(M'\) ，自然同态 \(\operatorname{Hom}_{\mathbf{A}}(\mathbf{M}, \mathbf{H}) \to \operatorname{Hom}_{\mathbf{A}}(\mathbf{M}', \mathbf{H})\) 是满射。
\end{enumerate}

（利用如下事实：对 A 的每个理想 b，存在一个整数 \(k \geqslant 0\) 使得 \(\mathfrak{b} \cap \mathfrak{a}^k \subset \mathfrak{ba}\) （AC，III，§ 3，第 1 号，推论 2）。）

\begin{enumerate}
\def\labelenumi{\alph{enumi})}
\setcounter{enumi}{1}
\item
  设 M 为一个内射 A-模。对于 \(n \geqslant 1\) ，我们用 \(\mathbf{H}_{n}\) 表示 M 中由被 \(\alpha^{n}\) 零化的元素组成的子模；令 \(\mathbf{H} = \bigcup \mathbf{H}_{n}\) 。证明 H 是一个内射 A-模。
\item
  证明 A/a 的内射包络的每个元素都被 a 的某个幂零化。* 30) 假设环 A 是局部、交换且诺特的；用 m 表示其极大理想，并令 k = A/m。设 I 为一个 A-模，使得 I 的每个元素都被 m 的某个幂零化。对每个 A-模 M，用 T(M) 表示 A-模 Hom\_A (M, I)，并用 c\_M : M → T(T(M)) 表示由
\end{enumerate}

\[
\left(c _ {\mathbf {M}} (m)\right) (u) = u (m) \quad \text { for } \quad m \in \mathbf {M}  , \quad u \in \mathrm{T} (\mathbf {M})  .
\]

定义的同态。证明下列性质是等价的：

α) 对每个有限长度 A-模 M，A-模 T(M) 是有限长度的，且同态 \(c_{M}: M \to T(T(M))\) 是双射。

β) 对每个有限长度的 A-模 M，有 long T(M) = long M。

γ) A-模 I 是内射的，且 A-模 k 与 T(k) 同构。

δ) A-模 I 与 k 的内射包络同构。

当这些条件满足时，有时称 I 为环 A 的一个对偶化模。（证明性质 γ) 与其他三个性质中的每一个等价：为看出 α) 或 β) 蕴含 γ)，使用习题 29。）

\begin{enumerate}
\def\labelenumi{\arabic{enumi})}
\setcounter{enumi}{30}
\item
  在上一习题的假设下，设 B 为一个局部、交换、诺特环，并设 \(\rho: A \to B\) 为一个环同态，使得 B 成为有限生成 A-模。若 I 是 A 的一个对偶化模，证明 B-模 \(\mathrm{Hom}_{\mathbf{A}}(\mathbf{B}, I)\) 是 B 的对偶化模。特别地，对 A 的每个理想 \(a\) ，I 中由被 \(a\) 零化的元素组成的子模是 A/a 的一个对偶化模。
\item
  \begin{enumerate}
  \def\labelenumii{\alph{enumii})}
  \tightlist
  \item
    我们保留习题 30 的假设；进一步假设 A 包含一个域 \(k_0\) 使得 \(k\) 是有限维的 \(k_0\) -向量空间。证明 A-模 \(l = \varinjlim_{n} \operatorname{Hom}_{k_0}(A/m^n, k_0)\)
  \end{enumerate}
\end{enumerate}

是 A 的对偶化模。

\begin{enumerate}
\def\labelenumi{\alph{enumi})}
\setcounter{enumi}{1}
\tightlist
\item
  设 \(k\) 为特征为零的域，B 为形式幂级数环 \(k[[\mathrm{T}_{1}, \ldots, \mathrm{T}_{n}]]\) 。我们赋予群 \(\mathbf{I} = k[\mathrm{X}_{1}, \ldots, \mathrm{X}_{n}]\) 一个 B-模结构，通过设定 \(\mathrm{T}_{i} \cdot \mathbf{P} = \partial \mathbf{P} / \partial \mathbf{X}_{i}\) 对所有 \(\mathbf{P} \in k[\mathrm{X}_{1}, \ldots, \mathrm{X}_{n}]\) 。证明 I 是 \(k\) 的一个内射包络。
\end{enumerate}

\exercisesection{A-模复形}

\begin{enumerate}
\def\labelenumi{\arabic{enumi})}
\tightlist
\item
  假设环 A 是主理想环；设 C 是自由 A-模的一个复形。
\end{enumerate}

\begin{enumerate}
\def\labelenumi{\alph{enumi})}
\item
  证明 C 是一族复形 \((^{p}\mathbf{E})_{p\in \mathbb{Z}}\) 的直和，满足： \(^{p}\mathbf{E}_{n} = 0\) 当 \(n\neq p\) , \(p - 1 ; Z_{n}(^{p}\mathbf{E}) = 0\) 。
\item
  我们进一步假设 \(C_n\) 对所有 \(n \in \mathbf{Z}\) 都是有限生成的。证明 \(C\) 是一些复形的直和，这些复形仅有一个非零分量或两个连续的非零分量，这些分量都同构于 \(A\) 。
\end{enumerate}

\begin{enumerate}
\def\labelenumi{\arabic{enumi})}
\setcounter{enumi}{1}
\item
  设 C 与 C' 为两个 A-模复形；假设复形 C 与 B(C)（相应地，C' 与 Z(C')）是投射的（相应地，内射的）。证明从 H(C) 到 H(C') 的每个复形的态射都由从 C 到 C' 的某个态射诱导。
\item
  设 C 为 A-模复形；将 C 视为 A(ε)-模（X，第 27 页）。我们用
\end{enumerate}

\[
i: \mathbf {A} \rightarrow \mathbf {A} (\varepsilon)
\]

表示典范单射。

证明以下条件是等价的：

\begin{enumerate}
\def\labelenumi{(\roman{enumi})}
\item
  C 是 i-投射的（X，第 170 页，习题 19）。
\item
  它是 i-内射的（同上）。
\item
  复形 C 同伦于零。
\item
  存在一个分次 A-模 M 以及分次 A(ε)-模的同构 A(ε) ⊗A M → C。
\item
  存在一个分次 A-模 N 以及分次 A(ε)-模的同构 C → Hom\_A(A(ε), N)。
\end{enumerate}

\begin{enumerate}
\def\labelenumi{\arabic{enumi})}
\setcounter{enumi}{3}
\tightlist
\item
  \begin{enumerate}
  \def\labelenumii{\alph{enumii})}
  \tightlist
  \item
    设 C 为右有界 A-复形，且 C 与 H(C) 均为投射的。证明 C 是分裂的。
  \end{enumerate}
\end{enumerate}

\begin{enumerate}
\def\labelenumi{\alph{enumi})}
\setcounter{enumi}{1}
\tightlist
\item
  在环 B = A(ε)（X，第 27 页）上，我们考虑由 C\_n = B （对所有 n）且 d(b) = εb （对 b ∈ B）定义的复形 C。复形 C 是自由的且同调为零，但并非分裂的。
\end{enumerate}

\begin{enumerate}
\def\labelenumi{\arabic{enumi})}
\setcounter{enumi}{4}
\tightlist
\item
  设 C 为一个复形，并设 s 为 C 的一个分裂， \(\delta: \mathbf{C} \to \mathbf{C}\) 为一个次数为 -1 的分次同态，使得 \((d + \delta) \circ (d + \delta) = 0\) 。用 \(\mathbf{C}'\) 表示复形 \((\mathbf{C}, d + \delta)\) 。假设从 C 到 C 的 A-线性映射 \(\theta\) （定义为 \(\theta = 1_C + \delta s\) ）是双射。
\end{enumerate}

\begin{enumerate}
\def\labelenumi{\alph{enumi})}
\item
  为了使 \(s\theta^{-1}\) 成为 \(C'\) 的一个分裂，必要且充分条件是 \(\delta(Z(C)) \subset \theta(B(C))\) 。
\item
  假设 a) 的条件成立。证明存在直和分解： \(\mathbf{C} = \mathbf{B}(\mathbf{C}) \oplus \mathbf{Ker}(ds) = \mathbf{B}(\mathbf{C}') \oplus \mathbf{Ker}(ds)\) 。若 p（相应地， \(p'\) ）是像为 \(\mathbf{Ker}(ds)\) 且核为 \(\mathbf{B}(\mathbf{C})\) （相应地， \(\mathbf{B}(\mathbf{C}')\) ）的投影算子，则有 \(p' = p\theta^{-1}\) 。
\item
  我们进一步假设从 C 到 C 的 A-线性映射 \(\lambda\) （定义为 \(\lambda = 1_{C} + s\delta\) ）是双射。证明我们有 \(Z(C') = \lambda^{-1}(Z(C))\) 与 \(C = Z(C) \oplus \operatorname{Im}(sd) = Z(C') \oplus \operatorname{Im}(sd)\) 。若 \(q\) （相应地， \(q'\) ）是像为 \(Z(C)\) （相应地， \(Z(C')\) ）且核为 \(\operatorname{Im}(sd)\) 的投影算子，则有 \(q' = \lambda^{-1} q\) 。
\end{enumerate}

\begin{enumerate}
\def\labelenumi{\arabic{enumi})}
\setcounter{enumi}{5}
\tightlist
\item
  我们沿用习题 5 的记号；进一步假设 A 是交换的，复形 C 与 \(\mathbf{C}^{\prime}/\mathbf{B}(\mathbf{C}^{\prime})\) 是平坦的，并且存在 A 的一个理想 m，满足 \(\delta(\mathbf{C}) \subset \mathfrak{m}\mathbf{C}\) 与 \(\bigcap_{n} (\mathbf{B}(\mathbf{C}) + \mathfrak{m}^{n}\mathbf{C}) = \mathbf{B}(\mathbf{C})\) （若 A 是带极大理想 m 的局部诺特环且 \(C_{n}\) 对每个 n 都是有限生成的 A-模，则最后一个条件总满足（参见 AC，III，§ 3，n° 3，命题 6））。然后证明 \(s\theta^{-1}\) 是 \(C^{\prime}\) 的一个分裂。
\end{enumerate}

（平坦性假设蕴含 \(\operatorname{Im}(d + \delta) \cap \mathfrak{m}^n\mathbf{C} = (d + \delta)\, (\mathfrak{m}^n\,\mathbf{C})\) 对所有 \(n\) 成立；若 \(x \in \mathbf{Z}(\mathbf{C})\) ，通过归纳构造元素 \(y_n \in \mathbf{C}\) , \(z_n \in \mathbf{Z}(\mathbf{C}) \cap \mathfrak{m}^n\mathbf{C}\) ，使得 \(\delta(x) = \theta d(y_n) + \delta(z_n)\) 。利用习题 5 a) 得出结论。）

\begin{enumerate}
\def\labelenumi{\arabic{enumi})}
\setcounter{enumi}{6}
\tightlist
\item
  设 C 与 C' 是两个复形， \(s_{1}(\text{resp. } s')\) 为 C（相应地，C'）的一个分裂， \(u : C' \to C\) 为一个复形的态射。我们定义 Con(u) 的一个 A-自同态 \(\sigma\) ，其分次次数为 -1，通过设定
\end{enumerate}

\[
\sigma (y ^ {\prime}, x) = \left(- s ^ {\prime} (y ^ {\prime}), s (x) - s u s ^ {\prime} (y ^ {\prime})\right).
\]

证明为了使 \(\sigma\) 成为 Con \((u)\) 的一个分裂，必要且充分的条件是我们有 H(u) = 0。

\begin{enumerate}
\def\labelenumi{\arabic{enumi})}
\setcounter{enumi}{7}
\tightlist
\item
  设 \(0 \to M' \xrightarrow{\mu} M \xrightarrow{\nu} M'' \to 0\) 为一个 A-模的正合列，它可以被视为一个序列
\end{enumerate}

复形的正合列（在次数 ≠ 0 时为零）。证明复形的态射 φ : Con(u) → M''\,（X，第 39 页）是同伦等价，当且仅当正合列 0 → M' → M → M'' → 0 是分裂的。9) 设 u : C' → C 是 A-复形的一个态射。记 v : C → Coker(u) 为商映射。考虑分次 A-同态

\[
\alpha : (\operatorname{Ker} u) (- 1) \rightarrow \operatorname{Con} (u) \quad \text { and } \quad \varphi : \operatorname{Con} (u) \rightarrow \operatorname{Coker} (u)
\]

由 \(\alpha(x') = (x', 0)\) , \(\varphi(x', x) = v(x)\) （对 \(x' \in C'\) , \(x \bullet C\) ）定义。证明 \(\alpha, \varphi\) 是复形的态射，并且它们产生一个正合列

\[
\dots \rightarrow \mathrm{H} _ {n - 1} (\text { Ker } u) \xrightarrow {\mathrm{H} (\alpha)} \mathrm{H} _ {n} (\text { Con } (u)) \xrightarrow {\mathrm{H} (\varphi)} \mathrm{H} _ {n} (\text { Coker } u) \xrightarrow {\delta} \mathrm{H} _ {n - 2} (\text { Ker } u) \rightarrow \dots
\]

其中 δ 是相对于这两个正合列的连接同态的复合，

\[
0 \rightarrow \operatorname{Ker} u \rightarrow C ^ {\prime} \rightarrow \operatorname{Im} u \rightarrow 0 \quad \text { and } \quad 0 \rightarrow \operatorname{Im} u \rightarrow C \rightarrow \operatorname{Coker} u \rightarrow 0.
\]

\begin{enumerate}
\def\labelenumi{\arabic{enumi})}
\setcounter{enumi}{9}
\tightlist
\item
  设 \(u: \mathbf{C} \to \mathbf{C}'\) 为 A-复形的一个态射， \(\mathbf{P}\) 为一个投射 A-模， \(n\) 为一个整数 \(\geqslant 0\) ，
\end{enumerate}

\[
h: \mathrm{P} \rightarrow \mathrm{H} _ {n} (\text { Con } (u))
\]

一个 A-同态。证明存在一个复形 \(\tilde{C}\) 、复形的一个正合列

\[
0 \to \mathbf {C} \xrightarrow {i} \tilde {\mathbf {C}} \to \mathbf {P} (- n) \to 0
\]

以及一个态射 \(\tilde{u}:\tilde{\mathbf{C}}\to \mathbf{C}'\) ，使得 \(\tilde{u}\circ i = u\) ，并且态射 \(\mathbf{C}(i):\mathrm{Con}(u)\to \mathrm{Con}(\tilde{u})\) （由 \(i\) 诱导）满足以下性质：

α) 映射 \(\mathrm{H}_{p}(\mathrm{C}(i)) : \mathrm{H}_{p}(\mathrm{Con}(u)) \to \mathrm{H}_{p}(\mathrm{Con}(\tilde{u}))\) 对 \(p \neq n, n + 1\) 是双射的；

β) 存在一个正合列

\[
0 \rightarrow \mathrm{H} _ {n + 1} (\operatorname{Con} (u)) \xrightarrow {\mathrm{H} _ {n + 1} (\mathbf {C} (i))} \mathrm{H} _ {n + 1} (\operatorname{Con} (\tilde {u})) \rightarrow \mathrm{P} \xrightarrow {h} \mathrm{H} _ {n} (\operatorname{Con} (u)) \xrightarrow {\mathrm{H} _ {n} (\mathbf {C} (i))} \mathrm{H} _ {n} (\operatorname{Con} (\tilde {u})) \rightarrow 0.
\]

\begin{enumerate}
\def\labelenumi{\arabic{enumi})}
\setcounter{enumi}{10}
\tightlist
\item
  我们把 A-模的双复形称为一个三元组 (C, \(d'\) , \(d''\) )，其中 C 是一个类型为 \(\mathbf{Z} \times \mathbf{Z}\) 的 A-分次模， \(d': \mathbf{C} \to \mathbf{C}\) 是一个次数为 (-1, 0) 的分次 A-自同态，且 \(d'': \mathbf{C} \to \mathbf{C}\) 是一个次数为 (0, -1) 的分次 A-自同态，满足 \(d' \circ d' = d'' \circ d'' = d' \circ d'' + d'' \circ d' = 0\) 。我们还将使用由 \(\mathbf{C}^{p,q} = \mathbf{C}_{-p,-q}\) （对 \(p, q \in \mathbf{Z}\) ）定义的上升分次。
\end{enumerate}

我们用 \(Z'(C)\) , \(Z''(C)\) , \(B'(C)\) , \(B''(C)\) 表示子双复形 \(\text{Ker } (d')\) , \(\text{Ker } (d'')\) , \(\text{Im } (d')\) , \(\text{Im } (d'')\) ，并用 \(H'(C)\) （相应地， \(H''(C)\) ）表示双复形 \(Z'(C)/B'(C)\) （相应地， \(Z''(C)/B''(C)\) ）。

\begin{enumerate}
\def\labelenumi{\alph{enumi})}
\tightlist
\item
  我们设 \(d = d' + d''\) 。证明当 C 装备上与给定双分次相关联的全分次时（由 \(\mathbf{C}_n = \bigoplus_{p+q=n} \mathbf{C}_{p,q}\) 定义，参见 II，第 164 页），二元组 (C, d) 是一个复形。它称为
\end{enumerate}

与双复形 (C, \(d'\) , \(d''\) ) 相关联的复形；其同调记为 H(C)。

\begin{enumerate}
\def\labelenumi{\alph{enumi})}
\setcounter{enumi}{1}
\item
  设 (C, \(d'\) , \(d''\) ) 与 (K, \(\delta'\) , \(\delta''\) ) 是两个双复形。从 (C, \(d'\) , \(d''\) ) 到 (K, \(\delta'\) , \(\delta''\) ) 的一个态射是从 C 到 K 的一个次数为 (0, 0) 的分次 A-同态 \(u\) ，满足： \(\delta'\circ u = u \circ d'\) 与 \(\delta''\circ u = u \circ d''\) 。证明 \(u\) 诱导了相应复形之间的一个态射。
\item
  设 \(u, v\) 为从 C 到 K 的两个双复形的态射。称连接 \(u\) 与 \(v\) 的同伦为从 C 到 K 的一对 A-同态 \((s', s'')\) ，其次数分别为 (1, 0) 和 (0, 1)，使得
\end{enumerate}

\[
g - f = d ^ {\prime} \circ s ^ {\prime} + s ^ {\prime} \circ d ^ {\prime} + d ^ {\prime \prime} \circ s ^ {\prime \prime} + s ^ {\prime \prime} \circ d ^ {\prime \prime}; \quad s ^ {\prime} \circ d ^ {\prime \prime} + d ^ {\prime \prime} \circ s ^ {\prime} = 0; \quad s ^ {\prime \prime} \circ d ^ {\prime} + d ^ {\prime} \circ s ^ {\prime \prime} = 0.
\]

证明 \(s' + s''\) 是连接由 \(u\) 与 \(v\) 诱导的相应复形的态射的同伦。12) 设 \((\mathbf{C}, d', d'')\) 为一个双复形，使得 \(\mathrm{H}_{p,-p}''(\mathbf{C}) = 0\) 对所有 \(p \in \mathbf{Z}\) 成立， \(\mathrm{H}_{-q,q+1}'(\mathbf{C}) = 0\) 当 \(q \geqslant 0\) ， \(\mathrm{H}_{r,-r-1}'(\mathbf{C}) = 0\) 当 \(r \geqslant 0\) ，并且存在正整数 \(a\) 与 \(b\) 使得 \(\mathbf{C}_{a,-a} = \mathbf{C}_{-b,b} = 0\) 。证明 A-模 \(\mathrm{H}_{0,0}'(\mathbf{C})\) 为零。

\begin{enumerate}
\def\labelenumi{\arabic{enumi})}
\setcounter{enumi}{12}
\tightlist
\item
  设 I 为一个有限集；我们用 \((e_i)_{i\in I}\) 表示 Z-模 \(\mathbf{Z}^{(I)}\) 的典范基。我们称一个 A-模的 I-复形（相应地，I-预复形）为一个类型为 \(\mathbf{Z}^{I}\) 的分次 A-模 C，配备一族 A-自同态 \((d_i)_{i\in I}\) ，使得 \(d_i\) 的分次次数为 \((-e_i)\) ，满足：
\end{enumerate}

\[
d _ {i} \circ d _ {i} = 0 \quad \text { for every } \quad i \in I;
\]

\[
d _ {i} \circ d _ {j} + d _ {j} \circ d _ {i} = 0 (\text { resp. } d _ {i} \circ d _ {j} = d _ {j} \circ d _ {i}) \quad \text { for } \quad i, j \in I.
\]

若 Card (I) = 1，则一个 I-复形就是一个复形；若 Card (I) = 2，则一个 I-复形就是一个双复形（习题 11）。若 \((\mathbf{C}, (d_i))\) 与 \((\mathbf{C}', (d_i'))\) 是两个 I-复形（相应地，I-预复形），则从 \((\mathbf{C}, (d_i))\) 到 \((\mathbf{C}', (d_i'))\) 的一个态射是一个次数为零的分次同态 \(u: \mathbf{C} \to \mathbf{C}'\) ，使得 \(d_i' \circ u = u \circ d_i\) 对所有 \(i \in I\) 成立。

\begin{enumerate}
\def\labelenumi{\alph{enumi})}
\tightlist
\item
  设 (C, \((d_i)\) ) 为一个 I-复形。若赋予 C 由其分次导出的全分次（Z 型），证明 (C, \(\sum_{i\in I}d_i\) ) 是一个 A-模复形，称为与 I-复形 (C, \((d_i)\) ) 相关联的复形。
\end{enumerate}

证明 I-复形的态射诱导出相应复形的态射。

\begin{enumerate}
\def\labelenumi{\alph{enumi})}
\setcounter{enumi}{1}
\tightlist
\item
  设 C 是一个 I-预复形，并设 R 是 I 上的一个全序关系，记为 ≤。对于 \(k = (k_i)_{i \in I}, x \in C_k\) 与 \(i \in I\) ，我们设
\end{enumerate}

\(\delta_{i}(x) = (-1)^{\sum_{j < i} k_{j}} d_{i}(x)\) 。证明 \((C, (\delta_{i}))\) 是一个 I-复形，记为 \(C_{R}\) 。

\begin{enumerate}
\def\labelenumi{\alph{enumi})}
\setcounter{enumi}{2}
\item
  设 \(\mathbf{R}_{\alpha}, \mathbf{R}_{\beta}\) 为 I 上的两个全序关系，分别记为 \(\leqslant_{\alpha}\) 与 \(\leqslant_{\beta}\) 。设 \(u_{\beta \alpha}\) 为 C 的自同态，由 \(u(x) = (-1)^{\varepsilon(k)} x\) （对 \(x \in \mathbf{C}_k\) ，其中 \(\varepsilon(k) = \sum_{\substack{i < _{\alpha} j \\ j < _{\beta} i}} k_i k_j\) ）定义。证明 \(u_{\beta \alpha}\) 是从 \(\mathbf{C}_{\mathbf{R}_{\alpha}}\) 到 \(\mathbf{C}_{\mathbf{R}_{\beta}}\) 的 I-复形的态射。
\item
  设 S 为 I 上全体全序关系组成的集；我们在 S 上赋予平凡预序（即满足如下条件的预序： \(R_{\alpha} \leqslant R_{\beta}\) ，无论 \(R_{\alpha}, R_{\beta} \in S\) ）。证明系统 \((\mathbf{C}_{R\alpha}, u_{\alpha\beta})\) 是相对于 S 的 I-复形的投影系；其极限 C 是一个 I-复形，称为与 I-预复形 C 相关联的 I-复形。证明 I-预复形的每个态射都诱导出相关联 I-复形之间的一个态射。
\end{enumerate}

A-模的一个谱序列定义为 A-复形的一个序列 \((\mathrm{E}_r, d_r)_{r \geqslant 2}\) 以及一族映射 \(\alpha_r: \mathrm{H}(\mathrm{E}_r) \to \mathrm{E}_{r+1}\) ，使得：

\begin{enumerate}
\def\labelenumi{(\roman{enumi})}
\tightlist
\item
  复形 \(E_r\) 的分次是与一个双分次相关联的全分次
\end{enumerate}

\[
\mathrm{E} _ {r} = \bigoplus_ {p, q \in \mathbf {Z}} \mathrm{E} _ {r} ^ {p, q} (r \geqslant 2);
\]

我们有 \(d_r(E_r^{pq}) \subset E_r^{p + r,q - r + 1}\) 。

\begin{enumerate}
\def\labelenumi{(\roman{enumi})}
\setcounter{enumi}{1}
\tightlist
\item
  映射 \(\alpha_{r}\) 是双分次 A-模之间的同构。
\end{enumerate}

我们还将使用 \(E_{r}\) 的下降双分次：记 \(E_{pq}^{r} = E_{r}^{-p, -q}\) 。

\begin{enumerate}
\def\labelenumi{\alph{enumi})}
\tightlist
\item
  这里我们用 \(Z_{r+1}(E_r)\) （相应地， \(B_{r+1}(E_r)\) ）表示双分次 A-模 Ker \(d_r\) （相应地，Im \(d_r\) ）。设
\end{enumerate}

\[
p _ {r}: \mathrm{E} _ {r} \rightarrow \mathrm{E} _ {r} / \mathrm{B} _ {r + 1} (\mathrm{E} _ {r})
\]

商映射以及 \(i_{r}: \mathbf{E}_{r+1} \to \mathbf{E}_{r}/\mathbf{B}_{r+1}(\mathbf{E}_{r})\) 为 \(\alpha_{r}^{-1}\) 与自然单射的复合映射。对每个 \(r \geqslant 2\) ，我们通过对 \(k\) 的归纳，如下定义双分次子模 \(\mathbf{Z}_{r+k}(\mathbf{E}_{r})\) 与 \(\mathbf{B}_{r+k}(\mathbf{E}_{r})\) （它们是 \(\mathbf{E}_{r}\) 的子模）：

\[
\mathbf {Z} _ {r + k} (\mathrm{E} _ {r}) = p _ {r} ^ {- 1} \left(i _ {r} \left(\mathbf {Z} _ {r + k} (\mathrm{E} _ {r + 1})\right)\right); \quad \mathbf {B} _ {r + k} (\mathrm{E} _ {r}) = p _ {r} ^ {- 1} \left(i _ {r} \left(\mathbf {B} _ {r + k} (\mathrm{E} _ {r + 1})\right)\right).
\]

证明我们有以下包含关系：

\[
0 \subset \mathrm{B} _ {r + 1} (\mathrm{E} _ {r}) \subset \mathrm{B} _ {r + 2} (\mathrm{E} _ {r}) \subset \dots \subset \mathrm{Z} _ {r + 2} (\mathrm{E} _ {r}) \subset \mathrm{Z} _ {r + 1} (\mathrm{E} _ {r}) \subset \mathrm{E} _ {r}
\]

并且 \(\mathbf{E}_k\) 同构于 \(\mathbf{Z}_k(\mathbf{E}_r) / \mathbf{B}_k(\mathbf{E}_r)\) （对 \(k\geqslant r + 1\) ）。

我们设 \(Z_{\infty}(E_2) = \bigcap_{r \geqslant 3} Z_r(E_2), B_{\infty}(E_2) = \bigcup_{r \geqslant 3} B_r(E_2)\) 与 \(E_{\infty} = Z_{\infty}(E_2)/B_{\infty}(E_2)\) 。我们称谱序列是正则的，如果递减序列 \((Z_r(E_2^{pq}))_{r \geqslant 2}\) 对每一对 \((p, q)\) 都是平稳的。

\begin{enumerate}
\def\labelenumi{\alph{enumi})}
\setcounter{enumi}{1}
\item
  设 \(E = (E_r, d_r, \alpha_r)_{r \geqslant 2}\) 与 \(E' = (E_r', d_r', \alpha_r')_{r \geqslant 2}\) 为两个 A-模的谱序列；从 E 到 \(E'\) 的一个态射是一列 A-复形的态射 \(u_r : E_r \to E_r'\) ，它与双分次相容，使得 \(\alpha_r' \circ H(u_r) = u_{r+1} \circ \alpha_r\) 对所有 \(r \geqslant 2\) 成立。证明对于每一个 \(r \geqslant 2\) 以及每个 \(k \geqslant r + 1\) , \(u_r\) 诱导出从 \(Z_k(E_r)\) （相应地， \(B_k(E_r)\) ）到 \(Z_k(E_r')\) （相应地， \(B_k(E_r')\) ）的一个同态。特别地， \(u_2\) 诱导出同态 \(Z_\infty(u_2) : Z_\infty(E_2) \to Z_\infty(E_2')\) , \(B_\infty(u_2) : B_\infty(E_2) \to B_\infty(E_2')\) 与 \(u_\infty : E_\infty \to E_\infty'\) ；所有这些同态都与双分次相容。
\item
  我们假设 \(u_2\) 是一个同构。证明此时 \(u_r\) 对所有 \(r \geqslant 2\) 以及 \(\mathbf{B}_{\infty}(u_2)\) 都是同构。此外，若序列 E 和 E' 是正则的，则 \(\mathbf{Z}_{\infty}(u_2)\) 与 \(u_{\infty}\) 都是同构。
\item
  设 \(G = \bigoplus_{n \in Z} G^n\) 为一个分次 A-模，配备分次子 A-模的一个递减序列 \((F^p G)_{p \in \mathbb{Z}}\)
\end{enumerate}

使得 \(\bigcap_{p} F^{p} G = 0\) 且 \(\bigcup_{p} F^{p} G = G\) 。我们说谱序列 E 逼近 \((G, (F^{p} G)_{p \in \mathbb{Z}})\)

（或简称为 G），如果存在同构 \(\beta^{pq}: E_{\infty}^{pq} \to (F^p G)^{p+q}/(F^{p+1} G)^{p+q}\) 对所有

\[
(p, q) \in \mathbf {Z} \times \mathbf {Z}.
\]

称谱序列 E 收敛于 G，如果它逼近 G 并且此外：

\begin{enumerate}
\def\labelenumi{(\roman{enumi})}
\item
  谱序列 E 是正则的；
\item
  对于每个整数 n，存在一个整数 p，使得 \((\mathbf{F}^{p}\mathbf{G})^{n}=0\) 。
\end{enumerate}

假设 E（相应地 E′）收敛于 G（相应地 G′）。设 \(u: \mathrm{E} \to \mathrm{E}'\) 为谱序列的一个态射， \(v: \mathrm{G} \to \mathrm{G}'\) 为一个次数为 0 的分次同态，使得 \(v(\mathrm{F}^p \mathrm{G}) \subset \mathrm{F}^p \mathrm{G}'\) 对所有 \(p\) 成立，并且诱导的同态 \(\bar{v}^{p,n}: (\mathrm{F}^p \mathrm{G})^n / (\mathrm{F}^{p+1} \mathrm{G})^n \to (\mathrm{F}^p \mathrm{G}')^n / (\mathrm{F}^{p+1} \mathrm{G}')^n\) 满足 \(\bar{v}^{p,p+q} \circ \beta^{p,q} = \beta'^{p,q} \circ u_{\infty}^{p,q}\) 对每一对 \((p, q)\) 成立。证明如果 \(u_2^{pq}\) 对每一对 \((p, q)\) 都是同构的，则 \(v\) 是一个同构。

\begin{enumerate}
\def\labelenumi{\arabic{enumi})}
\setcounter{enumi}{14}
\tightlist
\item
  设 E 是一个 A-模的谱序列，它逼近一个分次 A-模 G（习题 14）。
\end{enumerate}

\begin{enumerate}
\def\labelenumi{\alph{enumi})}
\tightlist
\item
  假设 \(E_{2}^{p,q}=0\) 当 p\textless0。证明映射 \(\alpha_{r}\) 对每个 q 诱导出单射的 A-同态：
\end{enumerate}

\[
\mathrm{E} _ {\infty} ^ {0, q} \rightarrow \dots \rightarrow \mathrm{E} _ {3} ^ {0, q} \rightarrow \mathrm{E} _ {2} ^ {0, q}
\]

由此，通过与 \((\beta^{0q})^{-1}\) 复合以及与到商映射的复合，得到一个 A-同态 \(\gamma^q: G^q \to E_2^{0,q}\) 。

\begin{enumerate}
\def\labelenumi{\alph{enumi})}
\setcounter{enumi}{1}
\tightlist
\item
  假设 \(E_{2}^{pq} = 0\) 当 q \textless{} 0。证明映射 \(\alpha_{r}^{-1}\) 诱导出满射同态：
\end{enumerate}

\[
\mathrm{E} _ {2} ^ {p, 0} \rightarrow \mathrm{E} _ {3} ^ {p, 0} \rightarrow \dots \rightarrow \mathrm{E} _ {\infty} ^ {p, 0}
\]

由此，通过与 \(\beta^{p,0}\) 复合，得到一个 A-同态 \(\delta^p:E_2^{p,0}\to G^p\) 。c) 我们假设 a) 和 b) 的假设均成立，即 \(E_{2}^{pq} = 0\) （若 p \textless{} 0 或 q \textless{} 0）。证明序列

\[
0 \rightarrow \mathrm{E} _ {2} ^ {1, 0} \xrightarrow {\delta^ {1}} \mathrm{G} ^ {1} \xrightarrow {\gamma^ {1}} \mathrm{E} _ {2} ^ {0, 1} \xrightarrow {d _ {2}} \mathrm{E} _ {2} ^ {2, 0} \xrightarrow {\delta^ {2}} \mathrm{G} ^ {2}
\]

是正合的。

\begin{enumerate}
\def\labelenumi{\alph{enumi})}
\setcounter{enumi}{3}
\item
  假设 \(\mathbf{E}_2^{pq} = 0\) （若 \(p > 0\) 或 \(q > 0\) ）；定义一个与 \(c\) ) 中类似的正合列。
\item
  假设存在整数 \(r, d\) ，满足 \(d \geqslant r \geqslant 2\) ，使得 \(\mathrm{E}_{r}^{pq} = 0\) （若 \(p \neq 0\) , \(d\) ）。定义一个正合列：
\end{enumerate}

\[
\dots \rightarrow \mathrm{E} _ {r} ^ {d, n - d} \rightarrow \mathrm{G} ^ {n} \rightarrow \mathrm{E} _ {r} ^ {0, n} \rightarrow \mathrm{E} _ {r} ^ {d, n + 1 - d} \rightarrow \mathrm{G} ^ {n + 1} \rightarrow \dots .
\]

若同样 \(E_{r}^{pq} = 0\) （当 \(q \neq 0\) , d），证明存在正合列：

\[
\dots \rightarrow \mathrm{E} _ {r} ^ {n, 0} \rightarrow \mathrm{G} ^ {n} \rightarrow \mathrm{E} _ {r} ^ {n - d, d} \rightarrow \mathrm{E} _ {r} ^ {n + 1, 0} \rightarrow \dots .
\]

\begin{enumerate}
\def\labelenumi{\alph{enumi})}
\setcounter{enumi}{5}
\tightlist
\item
  假设 \(\mathbf{E}_2^{pq} = 0\) （若 \(q \neq 0\) ）（相应地， \(p \neq 0\) ）。证明同态
\end{enumerate}

\[
\delta^ {n}: \mathrm{E} _ {2} ^ {n, 0} \rightarrow \mathrm{G} ^ {n} (\text { resp. } \gamma^ {n}: \mathrm{G} ^ {n} \rightarrow \mathrm{E} _ {2} ^ {0, n})
\]

对每个 n 都是双射。

¶ 16) 设 C 是一个 A-复形， \((F^{p}C)_{p\in\mathbf{Z}}\) 是 C 的子复形的一个递减序列。

\begin{enumerate}
\def\labelenumi{\alph{enumi})}
\tightlist
\item
  对于 \(p,q,r\in\mathbf{Z}\) ，设 \(M_{r}^{pq}\) 为 \((F^{p}C)^{p+q}\) 中满足 \(dx\in F^{p+r}C\) 的元素 x 组成的集合；对 \(r\geqslant0\) ，我们设 \(E_{r}^{pq}=M_{r}^{pq}/(d(M_{r-1}^{p-r+1,q+r-2})+M_{r}^{p+1,q-1})\) ，并用 \(d_{r}:E_{r}^{pq}\to E_{r}^{p+r,q-r+1}\) 表示由 \(d\) 通过取商得到的同态。证明对 \(r\geqslant0\) 存在同构 \(\alpha_{r}:H(E_{r})\to E_{r+1}\) ，使得序列 \((E_{r},d_{r})_{r\geqslant2}\) （配备映射 \(\alpha_{r}\) ）是一个谱序列。我们有
\end{enumerate}

\[
\mathrm{E} _ {0} ^ {p q} = (\mathrm{F} ^ {p} \mathrm{C}) ^ {p + q} / (\mathrm{F} ^ {p + 1} \mathrm{C}) ^ {p + q}.
\]

\begin{enumerate}
\def\labelenumi{\alph{enumi})}
\setcounter{enumi}{1}
\item
  我们假设 \(\bigcup_{p} \mathbf{F}^{p} \mathbf{C} = \mathbf{C}\) 并且对于每个 \(n\) ，存在一个整数 \(p\) 使得 \((\mathbf{F}^{p} \mathbf{C})^{n} = 0\) 。证明在 \(a)\) 中定义的谱序列收敛（习题 14）于分次 A-模 \(\mathrm{H}(\mathbf{C})\) ，它配备分次子模序列 \(\mathbf{F}^{p} \mathbf{H}(\mathbf{C}) = \operatorname{Im} \mathbf{H}(i_{p})\) ，其中 \(i_{p}: \mathbf{F}^{p} \mathbf{C} \to \mathbf{C}\) 是典范单射。
\item
  若 \(\mathbf{F}^0\bar{\mathbf{C}} = \mathbf{C}\) ，则 A-模 \(\mathbf{E}_r^{pq}\) 对 \(p < 0\) 为零；若 \((\mathbf{F}^p\mathbf{C})^n = 0\) （当 \(p > n\) ），则有 \(\mathbf{E}_r^{pq} = 0\) （对 \(q < 0\) ）。d) 设 \(\mathbf{C}'\) 为第二个复形， \((\mathbf{F}^p\mathbf{C}')_{p\in \mathbb{Z}}\) 为 \(\mathbf{C}'\) 的子复形的一个递减序列， \(\mathbf{E}'\) 为相应的谱序列。设 \(u:\mathbf{C}\to \mathbf{C}'\) 是复形的态射，使得 \(u(\mathbf{F}^p\mathbf{C})\subset \mathbf{F}^p\mathbf{C}'\) 对所有 \(p\) 成立。证明 \(u\) 诱导谱序列的一个态射 \(\tilde{u}:\mathbf{E}\to \mathbf{E}'\) 。
\item
  沿用 \(d\) ) 的记号，设 \(v: \mathbf{C} \to \mathbf{C}'\) 是第二个态射，使得 \(v(\mathbf{F}^p \mathbf{C}) \subset \mathbf{F}^p \mathbf{C}'\) 对所有 \(p\) 成立，并设 \(h\) 为联系 \(u\) 与 \(v\) 的一个同伦。设 \(k\) 为一个整数 \(\geqslant 0\) ，使得 \(h(\mathbf{F}^p \mathbf{C}) \subset \mathbf{F}^{p-k} \mathbf{C}'\) 对所有 \(p\) 成立。证明 \(\tilde{u}_r = \tilde{v}_r\) 当 \(r > k\) 。
\end{enumerate}

\begin{enumerate}
\def\labelenumi{\arabic{enumi})}
\setcounter{enumi}{16}
\tightlist
\item
  设 C 是一个双复形（X，第 174 页，习题 11）。
\end{enumerate}

\begin{enumerate}
\def\labelenumi{\alph{enumi})}
\tightlist
\item
  我们定义与 C 相关联的复形的子复形的两个递减序列 \(({}^{\prime}\mathbf{F}^p\mathbf{C})_{p\in \mathbf{Z}}\) 与 \(({}^{\prime\prime}\mathbf{F}^q\mathbf{C})_{q\in \mathbf{Z}}\) ，通过设定
\end{enumerate}

\[
\left(^ {\prime} \mathrm{F} ^ {p} \mathrm{C}\right) ^ {n} = \bigoplus_ {i \geqslant p} \mathrm{C} ^ {i, n - i} \quad \left(^ {\prime \prime} \mathrm{F} ^ {q} \mathrm{C}\right) ^ {n} = \bigoplus_ {j \geqslant q} \mathrm{C} ^ {n - j, j}.
\]

设′E 与″E 为相应的谱序列（习题 16，a））。证明对每一对 \((p, q)\) ，A-模 \({}^{\prime}E_2^{pq}\) （相应地， \({}^{\prime\prime}E_2^{pq}\) ）同构于 \(H^{\prime\prime p}(H^{\prime q}(C))\) （相应地， \(H^{\prime p}(H^{\prime\prime q}(C))\) ）。

\begin{enumerate}
\def\labelenumi{\alph{enumi})}
\setcounter{enumi}{1}
\item
  假设存在一个整数 \(n\) 使得 \(C^{pq} = 0\) （当 \(p \geqslant n\) 或 \(q \leqslant n\) ）。证明谱序列 ′E 收敛于 H(C)。如果同样存在一个整数 \(m\) 使得 \(C^{pq} = 0\) （当 \(p \leqslant m\) 或 \(q \geqslant m\) ），则谱序列 ″E 收敛于 H(C)。
\item
  我们假设 \(C^{pq} = 0\) （当 p \textless{} 0 或 q \textless{} 0）。证明同态 \(\delta^{p}: E_{2}^{p0} \to H^{p}(C)\) （习题 15）在取同调后，由复形的典范单射 ``Z\^{}\{0\}(C) → C'' 诱导而来。同样地，同态 ``E\_\{2\}\^{}\{p0\} → H\^{}\{p\}(C)'' 来自单射 ``Z\^{}\{0\}(C) → C''。
\item
  设 \(\tilde{C}\) 为第二个双复形，′E 与 ″E 为两个相应的谱序列。证明双复形的每个态射 \(u: C \to \tilde{C}\) 诱导谱序列的态射 ′u : ′E → ′E 与 ″u : ″E → ″E。证明如果两个态射 \(u\) , \(v\) （从 \(C\) 到 \(\tilde{C}\) ）是同伦的，则相应的态射 ′u 与 ′v（相应地，″u 与 ″v）相等（利用习题 16，e））。
\end{enumerate}

\begin{enumerate}
\def\labelenumi{\arabic{enumi})}
\setcounter{enumi}{17}
\tightlist
\item
  设 E 是一个谱序列（习题 14），使得 A-模 \(\mathbf{E}_2\) 具有有限长度。设 \(\mathcal{C}\) 是有限长度 A-模的同构类组成的集合；证明有 \(\chi_{\mathcal{G}}(\mathbf{E}_2) = \chi_{\mathcal{G}}(\mathbf{E}_r) = \chi_{\mathcal{G}}(\mathbf{E}_{\infty})\) 对所有 \(r \geqslant 2\) 成立（X，第 41 页）。如果此外 E 收敛于一个分次 A-模 G，则后者具有有限长度，并且有 \(\chi_{\mathcal{G}}(\mathbf{G}) = \chi_{\mathcal{G}}(\mathbf{E}_2)\) 。
\end{enumerate}

¶ 19) 我们称一个单纯复形为一个对 (K, F)，其中 K 是一个集合，F 是 K 的有限子集组成的一个集合，称为 K 的面，使得一个面的每个子集都是一个面，且 K 的每个点至少属于一个面。从一个单纯复形 (K, F) 到一个单纯复形 (L, G) 的单纯映射是从 K 到 L 的一个映射，它把 K 的每个面映到 L 的一个面。a) 设 n 是一个整数 ≥ 0。我们称单纯复形 K 的一个 n-单形为 K 的元素的序列 (α₀, \ldots, αₙ)，使得诸 αᵢ 属于 K 的同一个面。我们用 Sₙ(K, A) 表示由 n-单形（n ≥ 0）生成的自由 A-模，并对 n \textless{} 0 设 Sₙ(K, A) = 0。我们如下定义一个 A-同态 \(d_{n}:\mathbf{S}_{n}(\mathbf{K},\mathbf{A})\to \mathbf{S}_{n - 1}(\mathbf{K},\mathbf{A})\) ：设定 \(d_{n} = 0\) 当 \(n\leqslant 0\) ，且 \(d_{n}(\alpha_{0},\dots,\alpha_{n}) = \sum_{i = 0}^{n}(-1)^{i}(\alpha_{0},\dots,\theta_{i},\dots,\alpha_{n})\)

对所有 \(n\) -单形 \((\alpha_0, \ldots, \alpha_n)\) （字母上方的符号表示该字母应被省略）。证明 \(d_{n-1} \circ d_n = 0\) 对所有 \(n\) 成立，使得 \((S(K, A), d)\) 是 A-模的一个复形。以同样的方式定义 A-复形 \(C(K, A)\) ，通过设定 \(C^n(K, A) = \text{Hom}_Z(S_n(K, Z), A)\) , \(\delta^n = \text{Hom}(d_n, 1_A)\) （对所有 \(n\) ）。对 \(n \geqslant 0\) ，有时用 \(H_n(K, A)\) （相应地， \(H^n(K, A)\) ）表示 A-模 \(H_n(S(K, A))\) （相应地， \(H^n(C(K, A)))\) 。 b) 设 K, L 为两个单纯复形， \(f: K \to L\) 为一个单纯映射。证明 \(f\) 诱导复形的态射 \(S(f): S(K, A) \to S(L, A)\) 与 \(C(f): C(L, A) \to C(K, A)\) ，从而诱导 A-同态 \(f_*: H(K, A) \to H(L, A)\) 与 \(f^*: H'(L, A) \to H'(K, A)\) 。

\begin{enumerate}
\def\labelenumi{\alph{enumi})}
\setcounter{enumi}{2}
\tightlist
\item
  设 \(g: \mathbf{K} \to \mathbf{L}\) 为第二个单纯映射。我们说 \(f\) 与 \(g\) 是单纯同伦的，如果对每个面 \(\sigma\) （属于 \(\mathbf{K}\) ），集合 \(f(\sigma) \cup g(\sigma)\) 是 \(\mathbf{L}\) 的一个面。证明态射 \(\mathbf{S}(f)\) 与 \(\mathbf{S}(g)\) （相应地， \(\mathbf{C}(f)\) 与 \(\mathbf{C}(g)\) ）此时是同伦的（通过设定 \(h\) 来定义一个同伦
\end{enumerate}

\[
h (\alpha_ {0}, \dots , \alpha_ {n}) = \sum_ {i = 0} ^ {n} (- 1) ^ {i} (f (\alpha_ {0}), \dots , f (\alpha_ {i}), g (\alpha_ {i}), \dots , g (\alpha_ {n}))
\]

对每个 n-单形 \((\alpha_{0}, \ldots, \alpha_{n})\) 。

\begin{enumerate}
\def\labelenumi{\alph{enumi})}
\setcounter{enumi}{3}
\tightlist
\item
我们称单纯复形 K 为锥形的，如果存在 K 中的一点 s，使得对 K 的每个面 σ，σ ∪ \{s\} 是一个面。证明此时存在从 S(K, A) 和从 C(K, A) 到 A 上的同伦等价。e) 设 D(K, A) 表示 S(K, A) 的子复形，它由使得两个 αᵢ 相等的单形 (α₀, \ldots, αₙ) 以及形如 (α₀, \ldots, αₙ) - ε(σ). (ασ(0), \ldots, ασ(n)) 的元素（对每个 n-单形 (α₀, \ldots, αₙ) 以及 \{0, \ldots, n\} 的每个置换 σ）生成。证明 D(K, A) 是 S(K, A) 的一个子复形，并且商映射 p : S(K, A) → S(K, A)/D(K, A) 是同伦等价（通过在 K 上选取一个全序，并把单形 (α₀, \ldots..., αₙ) 模 D(K, A) 的类对应到单形 ε(σ). (ασ(0), \ldots, ασ(n))—其中 σ 是满足 ασ(0) \textless{} \ldots{} \textless{} ασ(n) 的置换，若诸 αᵢ 互不相同；否则对应到 0—即可得到一个同伦逆）。
\end{enumerate}

*20) 设 R 为一个完备离散赋值环（AC，VI，§ 3，第 6 节），其剩余域 k 的特征为 p \textgreater{} 0。假设 R 的分式域的特征为 0。设 v 表示 R 的规范化赋值，m 为 R 的极大理想，R* 为 R 的单位乘法群。对每个整数 \(n \geqslant 0\) ，设 R*(n) 为 R* 中由满足 \(v(x - 1) \geqslant n\) 的元素 x 构成的子群。 a) 设 r 为整数， \(r > v(p)/(p - 1)\) 。证明级数

\[
\exp (x) = \sum_ {i = 0} ^ {+ \infty} x ^ {i} / i!
\]

对 \(x \in \mathfrak{m}^r\) 收敛，并定义了从加法群 \(\mathfrak{m}^r\) 到群 \(\mathbf{R}^*(r)\) 上的一个同构。 b) 假设 \(k\) 是一个具有 \(q\) 个元素的有限域。设 \(n\) 为一个整数 \(\geqslant 1\) , \(\mu_n\) 为元素 \(x\) （ \(\mathbf{R}^*\) 中）满足 \(x^n = 1\) 所构成的子群。证明群 \(\mu_n\) 与 \(\mathbf{R}^*/(\mathbf{R}^*)^n\) 都是有限的，并且有

\[
\operatorname{Card} \left(\mathrm{R} ^ {*} / \mathrm{R} ^ {* n}\right) / \operatorname{Card} \left(\mu_ {n}\right) = q v (n).
\]

（设 \(\mathcal{C}\) 为有限 \(\mathbf{Z}\) -模的类组成的集合， \(\varphi : \mathrm{K}(\mathcal{C}) \to \mathbf{Q}^*\) 为如下定义的同态

\[
\varphi ([ \mathrm{G} ]) = \text { Card } (\mathrm{G}) \quad \text { for } \quad \mathrm{G} \in \mathscr {C},
\]

\(C^{(r)}\) 为满足如下条件的 Z-模复形： \(C_{p}^{(r)} = 0\) （当 \(p \neq 0, 1\) ）, \(C_{0}^{(r)} = C_{1}^{(r)} = R^{*}(r)\) , \(d_{1}(x) = x^{n}\) （当 \(x \in R^{*}(r)\) ）。证明 \(\varphi(H(C^{(r)}))\) 与 r 无关，然后使用 a)。

\begin{enumerate}
\def\labelenumi{\arabic{enumi})}
\setcounter{enumi}{20}
\tightlist
\item
  假设环 A 是交换环 k 上的一个交换代数；A 的 k-导子组成的 A-模 \(D_{k}(A)\) 于是被等同于 A-模 \(\Omega_{A/k}^{1}\) 的对偶，这在 \(\Omega_{A/k}\) 上定义了一个 \(\Lambda_{\mathbf{A}}(\mathbf{D}_{k}(\mathbf{A}))\) 左模结构（参见 III，第 165 页）。证明公式
\end{enumerate}

\[
\begin{array}{l} \left(\mathbf {X} _ {0} \wedge \dots \wedge \mathbf {X} _ {p}\right) \rfloor d \omega = \sum_ {i = 0} ^ {p} (- 1) ^ {i} \mathbf {X} _ {i} \left(\left(\mathbf {X} _ {0} \wedge \dots \wedge \hat {\mathbf {X}} _ {i} \wedge \dots \wedge \mathbf {X} _ {p}\right) \rfloor \omega\right) \\ \quad + \sum_ {0 \leqslant i <   j \leqslant p} (- 1) ^ {i + j} \left([ \mathbf {X} _ {i}, \mathbf {X} _ {j} ] \wedge \mathbf {X} _ {0} \wedge \dots \wedge \hat {\mathbf {X}} _ {i} \wedge \dots \wedge \hat {\mathbf {X}} _ {r} \wedge \dots \wedge \mathbf {X} _ {p}\right) \rfloor d \omega . \end{array}
\]

其中 \(\omega\in\Omega_{\mathbf{A}/k}^{p},\mathbf{X}_{0},\ldots,\mathbf{X}_{p}\in\mathbf{D}_{k}(\mathbf{A})\) ，而字母上的符号 \(^{\wedge}\) 表示该字母应被省略。（化归到 \(p=1\) 的情形。）

\begin{enumerate}
\def\labelenumi{\arabic{enumi})}
\setcounter{enumi}{21}
\tightlist
\item
  我们沿用前一习题的约定；若 M 是一个 A-模， \(\nabla\) 为 M 上的一个联络且若 \(X \in D_k(A)\) ，我们用 \(\nabla_X\) 表示 M 的 \(k\) -自同态 \((1_M \otimes X) \circ \nabla\) 。
\end{enumerate}

\begin{enumerate}
\def\labelenumi{\alph{enumi})}
\item
  设 \(\mathbf{M}_1, \mathbf{M}_2\) 为两个 A-模，分别配备联络 \(^1\nabla\) , \(^2\nabla\) ；证明可以如下定义 \(\nabla\) （ \(\mathbf{M}_1 \oplus \mathbf{M}_2\) （相应地， \(\mathbf{M}_1 \otimes_{\mathbf{A}} \mathbf{M}_2\) ，相应地， \(\operatorname{Hom}_{\mathbf{A}}(\mathbf{M}_1, \mathbf{M}_2)\) ）上的一个联络），通过设定： \(\nabla(m_1 + m_2) = ^1\nabla(m_1) + ^2\nabla(m_2)\) （当 \(m_1 \in \mathbf{M}_1\) , \(m_2 \in \mathbf{M}_2\) ）（相应地， \(\nabla_{\mathbf{X}}(m_1 \otimes m_2)) = ^1\nabla_{\mathbf{X}}(m_1) \otimes m_2 + m_1 \otimes ^2\nabla_{\mathbf{X}}(m_2)\) 当 \(\mathbf{X} \in \mathbf{D}_k(\mathbf{A})\) , \(m_1 \in \mathbf{M}_1\) 与 \(m_2 \in \mathbf{M}_2\) ，相应地， \(\nabla_{\mathbf{X}}(u) = ^2\nabla_{\mathbf{X}} \circ u - u \circ ^1\nabla_{\mathbf{X}}\) 当 \(\mathbf{X} \in \mathbf{D}_k(\mathbf{A})\) , \(u \in \operatorname{Hom}_{\mathbf{A}}(\mathbf{M}_1, \mathbf{M}_2)\) 。
\item
  设 M 为一个 A-模， \(\nabla: \mathbf{M} \to \mathbf{M} \otimes_{\mathbf{A}} \Omega_{\mathbf{A}/k}^{1}\) 为 M 上的一个联络。若 X, Y ∈ D \(_{k}\) (A)，我们用 R(X, Y) 表示 M 的如下 A-自同态：R(X, Y)(m) = \textless R(m), X ∧ Y\textgreater{} 对所有 m ∈ M 成立。证明公式：R(X, Y) = {[}∇X, ∇Y{]} - ∇{[}X,Y{]} （利用习题 21）。
\item
  假设 A-模 M 是投射的，从而曲率同态 R 被等同于 \(\operatorname{End}_{\mathbf{A}}(\mathbf{M}) \otimes_{\mathbf{A}} \Omega_{\mathbf{A}/k}^{2}\) 的一个元素。若给 \(\operatorname{End}_{\mathbf{A}}(\mathbf{M})\) 配备 \(\nabla\) （在 \(a\) ) 中定义的联络），证明有 \(\nabla^{2}\mathbf{R} = 0\) 。
\end{enumerate}

\exercisesection{分解}

\begin{enumerate}
\def\labelenumi{\arabic{enumi})}
\tightlist
\item
  设 \(a, b\) 为 A 的两个元素，使得 \(a\) （相应地， \(b\) ）的左零化子是理想 \(Ab\) （相应地， \(Aa\) ）。证明序列
\end{enumerate}

\[
\dots \longrightarrow \mathrm{A} _ {s} \xrightarrow {\delta_ {a}} \mathrm{A} _ {s} \xrightarrow {\delta_ {b}} \mathrm{A} _ {s} \xrightarrow {\delta_ {a}} \mathrm{A} _ {s} \longrightarrow \dots \longrightarrow \mathrm{A} _ {s} \xrightarrow {\delta_ {a}} \mathrm{A} _ {s} \longrightarrow \mathrm{A} / \mathrm{A} a \longrightarrow 0,
\]

其中 \(\delta_{a}\) （相应地， \(\delta_{b}\) ）表示右乘 \(a\) （相应地， \(b\) ），定义了 A-模 A/Aa 的一个自由分解。证明此构造适用于以下两种情形：

\(\alpha)\) 设 \(\mathbf{A}_0\) 为一个环；我们取 \(\mathrm{A} = \mathrm{A}_0(\varepsilon)\) （X，第 27 页）， \(a = b = \varepsilon\) 。

β）设G为有限循环群，σ为G的一个生成元，k为交换环。我们取

\[
\mathbf {A} = k ^ {(\mathbf {G})}, \quad a = 1 - e _ {\sigma}, \quad b = \sum_ {g \in \mathbf {G}} e _ {g}.
\]

\begin{enumerate}
\def\labelenumi{\arabic{enumi})}
\setcounter{enumi}{1}
\tightlist
\item
  我们假设环 A 是交换的。设 B 为一个 A-代数；对 \(n \geqslant 0\) ，我们用 \(\mathcal{A}^{n}(\mathbf{B})\) 表示 \((n + 1)\) 个都等于 B 的模在 A 上的张量积。我们设 \(\mathcal{A}^{n}(\mathbf{B}) = 0\) （当 \(n < 0\) ），并设 \(\mathcal{A}(\mathbf{B}) = \bigoplus \mathcal{A}^{n}(\mathbf{B})\) 。我们通过如下设定，定义 \(\mathcal{A}(\bar{\mathbf{B}})\) 的一个次数递增 +1 的分次 A-自同态：
\end{enumerate}

\[
d (b _ {0} \otimes \dots \otimes b _ {n}) = \sum_ {i = 0} ^ {n + 1} (- 1) ^ {i} b _ {0} \otimes \dots \otimes b _ {i - 1} \otimes 1 _ {\mathbf {B}} \otimes b _ {i} \otimes \dots \otimes b _ {n} \quad \text { for } b _ {0}, \dots , b _ {n} \in \mathbf {B}.
\]

\begin{enumerate}
\def\labelenumi{\alph{enumi})}
\tightlist
\item
  证明 \(d \circ d = 0\) ，使得 \((\mathcal{A}(\mathbf{B}), d)\) 是 A-模的一个复形。若 M 是 A-模，我们记 \(\mathcal{A}(\mathbf{B}, \mathbf{M})\) 为满足 \(\mathcal{A}^{n}(\mathbf{B}, \mathbf{M}) = \mathcal{A}^{n}(\mathbf{B}) \otimes_{\mathbf{A}} \mathbf{M}\) 的复形，其微分为 \(d \otimes 1_{\mathbf{M}}\) 。 b) 假设代数 B 是忠实平坦的 A-模（X，第 169 页，习题 9）。证明复形 \(\mathcal{A}(\mathbf{B}, \mathbf{M})\) 定义了 A-模 M 的一个右分解。（只需验证复形 \(\mathbf{B} \otimes_{\mathbf{A}} \mathcal{A}(\mathbf{B}, \mathbf{M})\) 定义了 \(\mathbf{B} \otimes_{\mathbf{A}} \mathbf{M}\) 的一个分解；构造这两个复形之间的一个 A-线性同伦等价。）
\end{enumerate}

\begin{enumerate}
\def\labelenumi{\arabic{enumi})}
\setcounter{enumi}{2}
\tightlist
\item
  设 B 是一个环，a 是 B 的一个双边理想，b 是 B 的一个包含 a 的左理想。记 A 为环 B/a，用 \(p: \mathbf{B} \to \mathbf{A}\) 表示商同态，用 \(b'\) 表示理想 \(p(\mathbf{b}) \subset \mathbf{A}\) 。a) 考虑典范嵌入的如下序列：
\end{enumerate}

\[
\dots \rightarrow a ^ {n} b \rightarrow a ^ {n} \rightarrow a ^ {n - 1} b \rightarrow a ^ {n - 1} \rightarrow \dots \rightarrow a \rightarrow b \rightarrow B.
\]

证明通过与B/a作张量积得到的序列

\[
\dots \rightarrow \mathfrak {a} ^ {n} \mathfrak {b} / \mathfrak {a} ^ {n + 1} \mathfrak {b} \rightarrow \mathfrak {a} ^ {n} / \mathfrak {a} ^ {n + 1} \rightarrow \dots \rightarrow \mathfrak {a} / \mathfrak {a} ^ {2} \rightarrow \mathfrak {b} / \mathfrak {a b} \rightarrow \mathbf {A}
\]

定义了A-模A/b'的一个左分解。

\begin{enumerate}
\def\labelenumi{\alph{enumi})}
\setcounter{enumi}{1}
\item
  假设理想 a 和 b 是自由左 A-模。证明前述分解是 A-模 A/b' 的一个自由分解。
\item
  设 G 为一个群，k 为一个交换环， \(A = \dot{k}^{(G)}\) 为 G 在 k 上的代数， \(\varepsilon : A \to k\) 为 k-代数的同态，使得 \(\varepsilon(e_g) = 1\) 对所有 \(g \in G\) 成立。我们赋予 k 与同态 \(\varepsilon\) 相伴的 A-模结构（左模结构）。设 \((g_i)_{i \in I}\) 为 G 的一个生成族，F 为在 I 上构造的自由群， \(B = k^{(F)}\) , \(\pi : F \to G\) 为满足 \(\pi(i) = g_i\) （当 \(i \in I\) ）的同态， \(p : B \to A\) 为由 \(\pi\) 导出的 k-代数同态。证明理想 a = Ker(p) 与 b = Ker ( \(\varepsilon \circ p\) ) 是自由左 B-模；由此推导出 \(k^{(G)}\) -模 k 的一个自由分解。
\end{enumerate}

（根据 I，第 147 页，习题 20，群 Ker (π) 具有一个基本族 (rα)α∈J；证明元素 (eᵣα - 1)，其中 α ∈ J（相应地 (eᵢ - 1)，其中 i ∈ I）构成左 B-模 a（相应地 b）的一组基。）

\begin{enumerate}
\def\labelenumi{\alph{enumi})}
\setcounter{enumi}{3}
\tightlist
\item
  采用 c) 中的记号，我们设 \(\mathbf{R} = \operatorname{Ker}(\pi)\) ；令 \(\varphi: \mathbf{R} \to \mathbf{R}/(\mathbf{R}, \mathbf{R})\) 为商映射。我们赋予 \(k\) -模 \(k \otimes_{\mathbf{Z}} \mathbf{R}/(\mathbf{R}, \mathbf{R})\) 一个 A-模结构，使得
\end{enumerate}

\[
e _ {\pi (f)} \left(1 \otimes \varphi (r)\right) = 1 \otimes \varphi (f r f ^ {- 1})
\]

对所有 \(r \in \mathbb{R}\) , \(f \in \mathbb{F}\) 成立。证明存在 A-模的一个正合列：

\[
0 \rightarrow k \otimes_ {\mathbf {Z}} \mathrm{R} / (\mathrm{R}, \mathrm{R}) \rightarrow \mathrm{A} ^ {(I)} \rightarrow \mathrm{A} \rightarrow k \rightarrow 0.
\]

\begin{enumerate}
\def\labelenumi{\arabic{enumi})}
\setcounter{enumi}{3}
\tightlist
\item
  假设环 A 是交换的。
\end{enumerate}

设 G 为一个群，K 为一个单纯复形（X，第 177 页，习题 19）。G 在 K 上的作用是从 G 到由 K 到自身的双射单纯映射组成的群的一个同态。a) G 在 K 上的作用定义了 G 在 A-模 S(K, A) 上的作用（同上）；证明该作用使 S(K, A) 成为 A \(^{(G)}\) -模的一个复形。

\begin{enumerate}
\def\labelenumi{\alph{enumi})}
\setcounter{enumi}{1}
\item
  假设 G 在 K 的底集合上自由作用。证明 A \(^{(G)}\) -模 S \(_{n}\) (K, A) 是自由的。
\item
  假设单纯复形 K 是锥形的。证明 \(\mathbf{A}^{(\mathbf{G})}\) -复形 \(\mathbf{S}(\mathbf{K},\mathbf{A})\) 定义了 A 的一个左分解，其中 A 配备了 \(\mathbf{A}^{(\mathbf{G})}\) -模结构，使得 \(e_g a = a\) 对所有 \(g \in \mathbf{G}\) , \(a \in \mathbf{A}\) 成立。
\item
  设 \(|\mathbf{G}|\) 为单纯复形 \((\mathbf{G},\mathcal{F})\) ，其中 \(\mathcal{F}\) 是 G 的有限子集组成的集合。群 G 通过左平移作用在单纯复形 \(|\mathbf{G}|\) 上；证明 \(\mathbf{A}^{(\mathbf{G})}\) -复形 S(\textbar G\textbar, A) 是 \(\mathbf{A}^{(\mathbf{G})}\) -模 A 的一个自由分解。
\item
  设 \(\mathbf{B}(\mathbf{A}^{(G)},\mathbf{A})\) 为 \(\mathbf{A}^{(G)}\) -模 A 的标准分解；对 \(n\geqslant 0\) ， \(\mathbf{A}^{(G)}\) -模 \(\mathbf{B}_n(\mathbf{A}^{(G)},\mathbf{A})\) 等同于 \((\mathbf{A}^{(G)})^{\otimes (n + 1)}\) 。证明 A-线性同态 \(u:\mathrm{S}(|\mathrm{G}|,\mathrm{A})\to \mathrm{B}(\mathrm{A}^{(G)},\mathrm{A})\) （定义为
\end{enumerate}

\[
u \left(g _ {0}, \dots , g _ {n}\right) = g _ {0} \otimes g _ {0} ^ {- 1} g _ {1} \otimes \dots \otimes g _ {n - 1} ^ {- 1} g _ {n} \quad \text { for } g _ {0}, \dots , g _ {n} \text { in } G,
\]

）是 \(A^{(G)}\) -复形之间的一个同构。

\begin{enumerate}
\def\labelenumi{\arabic{enumi})}
\setcounter{enumi}{4}
\tightlist
\item
  \begin{enumerate}
  \def\labelenumii{\alph{enumii})}
  \tightlist
  \item
    设 \(0 \to K \to P \to M \to 0\) , \(0 \to K' \to P' \to M \to 0\) 为两个 A-模的正合列，其中 \(P\) 与 \(P'\) 是投射的。证明存在一个交换图：
  \end{enumerate}
\end{enumerate}

\[
\begin{array}{c c c}0 \rightarrow \mathrm{K} \oplus \mathrm{P} ^ {\prime} \rightarrow \mathrm{P} \oplus \mathrm{P} ^ {\prime} \rightarrow \mathrm{M} \rightarrow 0\\\downarrow^ {v}&\downarrow^ {u}&\downarrow^ {l _ {M}}\\0 \rightarrow \mathrm{P} \oplus \mathrm{K} ^ {\prime} \rightarrow \mathrm{P} \oplus \mathrm{P} ^ {\prime} \rightarrow \mathrm{M} \rightarrow 0\end{array}
\]

其中水平行是正合的，并且 \(u\) 与 \(v\) 是同构（参见 II，第 183 页，习题 4）。b) 设 \(n\) 为一个整数 \(\geqslant 0\) , \((\mathbf{P}, p)\) 与 \((\mathbf{P}', p')\) 为 A-模 M 的两个左分解，使得 \(\mathbf{P}_j = \mathbf{P}_j' = 0\) （当 \(j > n\) ）且 A-模 \(\mathbf{P}_i\) 与 \(\mathbf{P}_i'\) 对 \(0 \leqslant i \leqslant n - 1\) 是投射的。证明 A-模 \(\bigoplus_{i \geqslant 0} (\mathbf{P}_{2i} \oplus \mathbf{P}_{2i+1}')\) 与 \(\bigoplus_{i \geqslant 0} (\mathbf{P}_{2j}^{\prime} \oplus \mathbf{P}_{2j+1})\) 是同构的。

\begin{enumerate}
\def\labelenumi{\alph{enumi})}
\setcounter{enumi}{2}
\tightlist
\item
  在 \(b\) ) 的假设下，证明存在两个零同调的投射复形 \(\mathbf{R}\) , \(\mathbf{R}'\) ，使得 \(\mathbf{R}_i = \mathbf{R}_i' = 0\) （当 \(i \notin [0, n]\) ），以及分解之间的一个同构 \(u: \mathbf{P} \oplus \mathbf{R} \to \mathbf{P}' \oplus \mathbf{R}'\) 。
\end{enumerate}

¶6). 设 M 为一个左 A-模。称长度为 n 的表示，或 M 的 n-表示，为一个正合列

\[
\mathrm{L} _ {n} \rightarrow \mathrm{L} _ {n - 1} \rightarrow \dots \rightarrow \mathrm{L} _ {1} \rightarrow \mathrm{L} _ {0} \rightarrow \mathrm{M} \rightarrow 0
\]

其中 \(L_{i}\) 是自由左 A-模（ \(0 \leqslant i \leqslant n\) ）。称该表示为有限的，如果所有 \(L_{i}\) 都是有限生成的自由模。

如果 M 是一个有限生成的左 A-模，则用 \(\lambda(M)\) （有限或等于 \(+\infty\) ）表示整数 \(n \geqslant 0\) （M 对之具有有限 \(n\) -表示）的上确界。如果 M 不是有限生成的，则令 \(\lambda(M) = -1\) 。

\begin{enumerate}
\def\labelenumi{\alph{enumi})}
\item
  设 \(0 \to \mathbf{P} \to \mathbf{N} \to \mathbf{M} \to 0\) 为左 A-模的一个正合列。则有 \(\lambda(\mathbf{N}) \geqslant \inf (\lambda(\mathbf{P}), \lambda(\mathbf{M}))\) 。
\item
  设 \(L_{n} \xrightarrow{d_{n}} L_{n-1} \to \cdots \to L_{0} \to M \to 0\) 为 M 的一个有限 n-表示；如果 \(\lambda(M) > n\) ，证明 Ker \((d_{n})\) 是有限生成的 A-模（利用习题 5，b））。
\item
  证明：在 a) 的假设下，有
\end{enumerate}

\[
\lambda (\mathbf {M}) \geqslant \inf \left(\lambda (\mathbf {N}), \lambda (\mathbf {P}) + 1\right).
\]

\((\mathrm{If} n \leqslant \inf (\lambda(\mathbf{N}), \lambda(\mathbf{P}) + 1)\) ，通过对 \(n\) 作归纳证明 \(\lambda(\mathbf{M}) \geqslant n\) ，推理方式如同在 \(a\) 中并使用 \(b)\) 。）

\begin{enumerate}
\def\labelenumi{\alph{enumi})}
\setcounter{enumi}{3}
\tightlist
\item
  证明：在 a) 的假设下，有
\end{enumerate}

\[
\lambda (\mathbf {P}) \geqslant \inf \left(\lambda (\mathbf {N}), \lambda (\mathbf {M}) - 1\right).
\]

（方法与 c）相同。）由此推出，若 \(\lambda(\mathbf{N}) = +\infty\) ，则 \(\lambda(\mathbf{M}) = \lambda(\mathbf{P}) + 1\) 。

\begin{enumerate}
\def\labelenumi{\alph{enumi})}
\setcounter{enumi}{4}
\tightlist
\item
  由 a)、c) 和 d) 推出，若 \(\mathbf{N} = \mathbf{M} \oplus \mathbf{P}\) ，则有
\end{enumerate}

\[
\lambda (\mathbf {N}) = \inf \left(\lambda (\mathbf {M}), \lambda (\mathbf {P})\right).
\]

特别地，为使 N 具有有限表示，必要且充分的是 M 与 P 也具有有限表示。f) 设 \(N_{1}\) , \(N_{2}\) 为 A-模 M 的两个子模。假设 \(N_{1}\) 与 \(N_{2}\) 都具有有限表示。为使 \(N_{1} + N_{2}\) 具有有限表示，必要且充分条件是 \(N_{1} \cap N_{2}\) 是有限生成的。

\begin{enumerate}
\def\labelenumi{\arabic{enumi})}
\setcounter{enumi}{6}
\tightlist
\item
  \begin{enumerate}
  \def\labelenumii{\alph{enumii})}
  \tightlist
  \item
    使用习题 6 的记号，证明：若 M 是投射模，则有
  \end{enumerate}
\end{enumerate}

\[
\lambda (\mathbf {M}) = - 1 \quad \text { or } \quad \lambda (\mathbf {M}) = + \infty .
\]

若 A 是左诺特环，则对每个 A-模 M，有 \(\lambda(M) = -1\) 或 \(\lambda(M) = +\infty\) 。 b) 若 a 是环 A 的一个左理想，且不是有限生成的，则 \(A_s/a\) 是一个单演 A-模，它不具有有限表示，换句话说 \(\lambda(A_s/a) = 0\) （参见 X，第 11 页，命题 6）。

\begin{enumerate}
\def\labelenumi{\alph{enumi})}
\setcounter{enumi}{2}
\item
  给出环 A 的一个单演左理想 a 的例子，使得 \(A_{s}/a\) （它是有限表示的）容许一个对偶，该对偶不是有限生成的右 A-模。
\item
  设 K 为一个交换域，E 为向量空间 \(\mathbf{K}^{(\mathbf{N})}\) , \((e_n)\) 为 E 的典范基，T 为 E 的张量代数，因此它的一组基由有限乘积 \(e_{i_1} e_{i_2} \ldots e_{i_k}\) 构成（ \(k \geqslant 0\) , \(i_j \in \mathbf{N}\) 对所有 \(j\) ）。对给定的整数 \(n\) ，设 b 为 T 的由乘积 \(e_1 e_0, e_2 e_1, \ldots, e_n e_{n-1}\) 与 \(e_{n+k} e_n\) （对所有 \(k \geqslant 1\) ）生成的双边理想；设 A 为商环 T/b，并对每个整数 \(m\) ，设 \(a_{m}\) 为 \(e_{m}\) 在 A 中的典范像。证明，若 \(\mathbf{M} = \mathbf{A}_{s}/\mathbf{A}a_{0}\) ，则 \(\lambda(\mathbf{M}) = n\) （注意，对 \(m \leqslant n - 1\) ， \(a_{m}\) 的左零化子是 \(\mathbf{A}a_{m+1}\) ，并利用习题 6，b)）。
\end{enumerate}

\begin{enumerate}
\def\labelenumi{\arabic{enumi})}
\setcounter{enumi}{7}
\tightlist
\item
  设 C 为一个交换环，E、F 为两个 C-模。证明有 \(\lambda(E \otimes_{C} F) \geqslant \inf(\lambda(E), \lambda(F))\) 。
\item
  设 \(\rho: A \to B\) 为一个环同态，M 是一个 A-模。
\end{enumerate}

\begin{enumerate}
\def\labelenumi{\alph{enumi})}
\item
  如果 B 是平坦的右 A-模且 M 具有有限 n-表示，则 B-模 B \(\otimes_{\mathbf{A}}\) M 具有有限 n-表示。
\item
  如果 B 是忠实平坦的右 A-模（X，第 169 页，习题 9），并且 B-模 B ⊗\_A M 具有有限 n-表示，则 M 具有有限 n-表示（利用习题 6，b）和第 169 页习题 12）。
\end{enumerate}

\begin{enumerate}
\def\labelenumi{\arabic{enumi})}
\setcounter{enumi}{9}
\tightlist
\item
  设 E 是一个 A-模。称 E 是伪凝聚的，如果 E 的每个有限生成子模都是有限表示的；伪凝聚模的每个子模都是伪凝聚的。称 E 是凝聚的，如果它是伪凝聚的且有限生成的（因此是有限表示的）。
\end{enumerate}

\begin{enumerate}
\def\labelenumi{\alph{enumi})}
\item
  设 \(0 \to E' \to E \to E'' \to 0\) 为 A-模的一个正合列。证明，如果 E 是伪凝聚的（相应地，凝聚的）且 \(E'\) 是有限生成的，则 \(E''\) 是伪凝聚的（相应地，凝聚的）。证明，如果 \(E'\) 与 \(E''\) 都是伪凝聚的（相应地，凝聚的），则 E 亦然。证明，如果 E 与 \(E''\) 是凝聚的，则 \(E'\) 亦然（使用 X，第 11 页，命题 6）。
\item
  设 E 是一个凝聚 A-模，E' 是一个伪凝聚（相应地，凝聚）A-模。证明：对每个同态 \(u: \mathrm{E} \to \mathrm{E}'\) , \(\operatorname{Im}(u)\) 与 \(\operatorname{Ker}(u)\) 都是凝聚的，且 Coker（ \(u\) ）是伪凝聚的（相应地，凝聚的）（使用 \(a\) )）。
\item
  证明：伪凝聚（相应地，凝聚）模的任意直和（相应地，任意有限直和）是伪凝聚（相应地，凝聚）模。
\item
  若 E 是伪凝聚模，且 M, N 是 E 的凝聚子模，证明 \(M + N\) 与 \(M \cap N\) 是凝聚的（使用 a) 和 c)）。
\item
  设 A 是交换的。证明：若 E 是凝聚 A-模，F 是凝聚（相应地，伪凝聚）A-模，则 Hom\_A(E, F) 是凝聚（相应地，伪凝聚）A-模。（化归到 F 是凝聚的情形；考虑 E 的一个有限表示，然后使用 b)。）
\end{enumerate}

¶11) a) 设 A 是一个环。证明下列性质是等价的：

\(\alpha)\) A-模 \(\mathbf{A}_s\) 是凝聚的（习题 10）。

β) 每个有限表示的 A-模都是凝聚的。

γ) 对每个有限表示的右 A-模 M，对偶 M* 是有限生成的 A-模。

δ) 对每个集合 I，右 A-模 \(A_{d}^{I}\) 是平坦的。

ε) 平坦右 A-模的任意乘积是平坦的。

（为了证明 \(\alpha\) ）蕴含 \(\beta\) ），使用习题 10，b)。为了证明 \(\delta\) ）蕴含 \(\gamma\) ），使用 X，第 14 页，定理 1；为了证明 \(\alpha\) ）蕴含 \(\varepsilon\) ），使用 X，第 170 页，习题 18。）

满足这些性质的环称为左凝聚的（或简称为凝聚的），右凝聚环的概念类似地定义。

\begin{enumerate}
\def\labelenumi{\alph{enumi})}
\setcounter{enumi}{1}
\item
  证明左诺特环是凝聚的。给出一个不是凝聚环的右阿廷环的例子（参见 VIII，§ 1，习题 4）。
\item
  证明绝对平坦环（X，第 169 页，习题 16）是凝聚的。
\item
  证明，若 A 是凝聚环且 M 是 A-模，则我们有 \(\lambda(M) = -1\) 或 \(\lambda(M) = 0\) 或 \(\lambda(M) = +\infty\) （习题 6）。特别地，每个有限表示的 A-模都容许一个由有限生成自由模构成的分解。
\item
  设 \((\mathbf{A}_{\alpha}, \varphi_{\beta \alpha})\) 为一个环的归纳系统，其指标集是滤过的，并设 \(\mathbf{A} = \varinjlim \mathbf{A}_{\alpha}\) 。我们假设对 \(\alpha \leqslant \beta\) , \(\mathbf{A}_{\beta}\) 是一个平坦的右 \(\mathbf{A}_{\alpha}\) -模。证明若诸 \(\mathbf{A}_{\alpha}\) 是凝聚的，则 \(\mathbf{A}\) 亦然（注意 \(\mathbf{A}\) 是一个平坦的右 \(\mathbf{A}_{\alpha}\) -模，对所有 \(\alpha\) 成立，且对每个有限生成的左理想 \(\alpha \subset \mathbf{A}\) ，存在一个指标 \(\alpha\) 与一个理想 \(\alpha_{\alpha}\) （ \(\mathbf{A}\) 中的），使得 \(\alpha\) 同构于 \(\mathbf{A} \otimes_{\mathbf{A}} a_{\alpha'}\) 。）
\item
  从 e) 推出，交换诺特环上的任意多项式环（对任意有限或无限个未定元）是凝聚的。由此推出，凝聚环的商环不一定是凝聚的。
\item
  为使 A 是左凝聚的，必要且充分条件是 A 的每个元素的左零化子是有限生成的，并且两个有限生成的左理想的交是有限生成的（使用习题 6，f））。
\end{enumerate}

\begin{enumerate}
\def\labelenumi{\arabic{enumi})}
\setcounter{enumi}{11}
\item
  设 c 为无穷基数。假设环 A 的每个左理想都容许一个基数 ≤ c 的生成集。证明：由基数 ≤ c 的元素集生成的每个 A-模都容许一个由秩 ≤ c 的自由 A-模构成的左分解（参见 VIII，§ 1，习题 13）。
\item
  设 M 是一个 A-模。M 的一个极小内射分解是 M 的一个内射分解 \((e, 1)\) ，使得 \(I^n\) 是 \(Z^n(I)\) 的内射包络（对所有 \(n \geqslant 0\) ）。
\end{enumerate}

\begin{enumerate}
\def\labelenumi{\alph{enumi})}
\tightlist
\item
  每个 A-模都具有极小内射分解；任意两个这样的分解都是同构的。
\item
  设 I、I' 为 M 的两个内射分解，其中 I 是极小的；设 \(f: I \to I'\) 与 \(g: I' \to I\) 为两个分解的态射。则 f 是单射，g 是满射，且它是子复形 Im(f)（同构于 I）与 Ker(g)（带零同调）的直和。
\end{enumerate}

\begin{enumerate}
\def\labelenumi{\arabic{enumi})}
\setcounter{enumi}{13}
\tightlist
\item
  假设环 A 是局部诺特环；用 m 表示其极大理想。设 L 是由有限生成 A-模组成的复形，它在次数 \textless{} 0 时为零，M 是一个有限生成的 A-模，且 q : L → M 是一个态射。我们作如下假设：
\end{enumerate}

\begin{enumerate}
\def\labelenumi{(\roman{enumi})}
\item
  同态 \(1 \otimes q : (\mathbf{A}/\mathfrak{m}) \otimes_{\mathbf{A}} \mathbf{L}_{0} \to (\mathbf{A}/\mathfrak{m}) \otimes_{\mathbf{A}} \mathbf{M}\) 是单射。
\item
  我们有 \(d_{i}(\mathbf{L}_{i})\subset \mathfrak{m}\mathbf{L}_{i - 1}\) 当 \(i\geqslant 1\) ，并且映射
\end{enumerate}

\[
\bar {d} _ {i}: \mathrm{L} _ {i} / \mathrm{mL} _ {i} \rightarrow \mathrm{mL} _ {i - 1} / \mathrm{m} ^ {2} \mathrm{L} _ {i - 1} \quad \text { is injective }.
\]

设 (P, p) 是 M 的一个投射分解， \(u : L \to P\) 为一个态射，使得 \(p \circ u = q\) 。证明 u 是单射，并且 \(u(L)\) 是 P 作为分次 A-模的一个直和项（化归到分解 P 为极小的情形，并利用 VIII，§ 8，第 3 号，推论 3）。

15) 设 A 是一个诺特局部环，m 为其极大理想，k = A/m。设 P 为 A-模 k 的一个极小投射分解；对 \(n \geqslant 0\) ，我们用 \(b_{n}\) 表示自由 A-模 \(P_{n}\) 的秩。

\begin{enumerate}
\def\labelenumi{\alph{enumi})}
\item
  证明我们有 \(b_0 = 1\) ，并且 \(b_1 = \dim_k (m/m^2)\) 是 \(m\) 的生成元的最小数目。
\item
  若 A 是交换的，我们有 \(b_{2} \geqslant \frac{1}{2} b_{1}(b_{1} - 1)\) （考虑正合列 \(\mathbf{P}_{2} \to m\mathbf{P}_{1} \to m^{2} \to 0\) ）。 c) 假设环 A 是阿廷环（不一定是交换的）。证明不等式 \(b_{2} \geqslant \frac{1}{4} b_{1}^{2}\) （设 \(\mathbf{M}_{l}\) 为 \(\mathbf{P}_{2}\) 的由元素 \(x\) （满足 \(d_{2}(x) \in m^{l}\mathbf{P}_{1}\) ）组成的子模；证明我们有
\end{enumerate}

\[
\mathbf {P} _ {2} \subset \mathbf {M} _ {l + 1}
\]

并且存在一个正合列 \(0 \to \mathbf{M}_l / \mathbf{M}_{l+1} \to \mathfrak{m}^l \mathbf{P}_1 / \mathfrak{m}^{l+1} \mathbf{P}_1 \to \mathfrak{m}^{l+1} / \mathfrak{m}^{l+2} \to 0\) 。由此推出多项式 \(p(t) = \sum_{l} (\dim_k (\mathfrak{m}^l / \mathfrak{m}^{l+1})) t^l\) 满足 \((b_2 t^2 - b_1 t + 1)p(t) \geqslant 0\) （当 \(0 \leqslant t \leqslant 1\) ）。通过分别考察情形

 \(b_{1} = 1, b_{1} \geqslant 2\) ，得出结论。

\begin{enumerate}
\def\labelenumi{\arabic{enumi})}
\setcounter{enumi}{15}
\tightlist
\item
  \begin{enumerate}
  \def\labelenumii{\alph{enumii})}
  \tightlist
  \item
    设 \(0 \to \mathbf{M}' \to \mathbf{M} \to \mathbf{M}'' \to 0\) 为 A-模的一个正合列， \((\mathbf{P}', p')\) （相应地， \((\mathbf{P}'', p'')\) ）为 \(\mathbf{M}'\) （相应地， \(\mathbf{M}''\) ）的一个投射分解。证明存在一个投射分解 \((\mathbf{P}, p)\) （ \(\mathbf{M}\) 的）以及一个交换图：
  \end{enumerate}
\end{enumerate}

\[
\begin{array}{c} 0 \to \mathrm{P} ^ {\prime} \to \mathrm{P} \to \mathrm{P} ^ {\prime \prime} \to 0 \\ \downarrow^ {p ^ {\prime}} \quad \downarrow^ {p} \quad \downarrow^ {p ^ {\prime \prime}} \\ 0 \to \mathrm{M} ^ {\prime} \to \mathrm{M} \to \mathrm{M} ^ {\prime \prime} \to 0 \end{array}
\]

其中水平行是正合的。

\begin{enumerate}
\def\labelenumi{\alph{enumi})}
\setcounter{enumi}{1}
\tightlist
\item
  设 \(0 \to \mathbf{N}' \to \mathbf{N} \to \mathbf{N}'' \to 0\) 为 A-模的第二个正合列， \((\mathbf{Q}, q)\) （相应地， \((\mathbf{Q}', q')\) ，相应地， \((\mathbf{Q}'', q'')\) ）为 \(\mathbf{N}\) （相应地， \(\mathbf{N}'\) ，相应地， \(\mathbf{N}''\) ）的一个投射分解，从而存在一个具有正合行的交换图：
\end{enumerate}

\[
\begin{array}{c}0 \rightarrow Q ^ {\prime} \rightarrow Q \rightarrow Q ^ {\prime \prime} \rightarrow 0\\\downarrow^ {q ^ {\prime}} \quad \downarrow^ {q} \quad \downarrow^ {q ^ {\prime \prime}}\\0 \rightarrow N ^ {\prime} \rightarrow N \rightarrow N ^ {\prime \prime} \rightarrow 0\end{array}
\]

另一方面，设：

\[
\begin{array}{c} 0 \to \mathrm{M} ^ {\prime} \to \mathrm{M} \to \mathrm{M} ^ {\prime \prime} \to 0 \\ \downarrow^ {u ^ {\prime}} \quad \downarrow^ {u} \quad \downarrow^ {u ^ {\prime \prime}} \\ 0 \to \mathrm{N} ^ {\prime} \to \mathrm{N} \to \mathrm{N} ^ {\prime \prime} \to 0. \end{array}
\]

它是正合列的一个交换图；设 \(\tilde{u}':\mathbf{P}'\to\mathbf{Q}'\) 与 \(\tilde{u}'':\mathbf{P}''\to\mathbf{Q}''\) 为复形的态射，使得 \(u'\circ p'=q'\circ\tilde{u}'\) 与 \(u''\circ p''=q''\circ\tilde{u}'\) 。证明存在一个态射 \(\tilde{u}:\mathbf{P}\to\mathbf{Q}\) ，使得 \(u\circ p=q\circ\tilde{u}\) ，并且下图：

\[
0 \rightarrow \mathrm{P} ^ {\prime} \rightarrow \mathrm{P} \rightarrow \mathrm{P} ^ {\prime \prime} \rightarrow 0\tag{*}
\]

\[
0 \rightarrow Q ^ {\prime} \rightarrow Q \rightarrow Q ^ {\prime \prime} \rightarrow 0
\]

是交换的。

\begin{enumerate}
\def\labelenumi{\alph{enumi})}
\setcounter{enumi}{2}
\tightlist
\item
  设 \(\tilde{u}_{1}^{\prime}: \mathrm{P}^{\prime} \to \mathrm{Q}^{\prime}, \tilde{u}_{1}: \mathrm{P} \to \mathrm{Q}\) 与 \(\tilde{u}_{1}^{\prime\prime}: \mathrm{P}^{\prime\prime} \to \mathrm{Q}^{\prime\prime}\) 为三个复形的态射，使得
\end{enumerate}

\[
u ^ {\prime} \circ p ^ {\prime} = q ^ {\prime} \circ \tilde {u} _ {1} ^ {\prime}, u \circ p = q \circ \tilde {u} _ {1}, u ^ {\prime \prime} \circ p ^ {\prime \prime} = q ^ {\prime \prime} \circ \tilde {u} _ {1} ^ {\prime \prime}
\]

使图 (*) 交换；设 \(s'\) （相应地， \(s''\) ）为连接 \(\tilde{u}'\) 与 \(\tilde{u}_{1}'\) （相应地， \(\tilde{u}''\) 与 \(\tilde{u}_{1}'\) ）的一个同伦。证明存在一个同伦 \(s\) 连接 \(\tilde{u}\) 与 \(\tilde{u}_{1}\) ，使得该图：

\[
\begin{array}{c}0 \rightarrow \mathrm{P} ^ {\prime} \rightarrow \mathrm{P} \rightarrow \mathrm{P} ^ {\prime \prime} \rightarrow 0\\\downarrow s ^ {\prime} \quad \downarrow s \quad \downarrow s ^ {\prime \prime}\\0 \rightarrow \mathrm{Q} ^ {\prime} \rightarrow \mathrm{Q} \rightarrow \mathrm{Q} ^ {\prime \prime} \rightarrow 0\end{array}
\]

是交换的。

\begin{enumerate}
\def\labelenumi{\alph{enumi})}
\setcounter{enumi}{3}
\tightlist
\item
  陈述并证明关于内射分解的与 a)、b)、c) 相对应的结果。 ¶17) 在本习题中，若 \((\mathbf{K}, d', d'')\) 是一个双复形（X，第 174 页，习题 11），我们将用 \(K_{p,*}\) 表示复形 \((\bigoplus_{n} K_{p,n}, d'')\) 。每个复形都将被视为一个双复形，它在次数 \((p, q)\) （当 \(q \neq 0\) ）处为零。
\end{enumerate}

设 C 为 A-模的一个复形。C 的一个投射分解是一个二元组 \((P, p)\) ，其中 P 是一个双复形且 \(p : P \to C\) 是双复形的一个态射，使得 \(P_{p,q} = 0\) （对 q \textless{} 0）且使得复形 \(P_{p,*}\) （相应地， \(Z'_{p,*}(\hat{P})\) , \(B'_{p,*}(\mathbf{P})\) , \(H'_{p,*}(\mathbf{P})\) ）定义了 \(C_p\) （相应地， \(Z_p(\mathbf{C})\) , \(B_p(\mathbf{C})\) , \(H_p(\mathbf{C})\) ）的一个投射分解（对每个 \(p \in Z\) ）。为使 \((P, p)\) 是 C 的一个投射分解，只需复形 \(B'_{p,*}(\mathbf{P})\) 与 \(H'_{p,*}(\mathbf{P})\) 对每个 p 定义 \(B_p(\mathbf{C})\) 与 \(H_p(\mathbf{C})\) 的投射分解即可。

\begin{enumerate}
\def\labelenumi{\alph{enumi})}
\tightlist
\item
  证明每个 A-复形 C 都容许一个投射分解（对每个 p，选取 \(B_{p,*}\) 与 \(H_{p,*}\) 为 \(B_{p}(C)\) 与 \(H_{p}(C)\) 的投射分解；利用习题 16，构造投射分解 \(Z_{p,*}\) （ \(Z_{p}(C)\) 的）与 \(P_{p,*}\) （ \(C_{p}\) 的），以及正合列：
\end{enumerate}

\[
0 \rightarrow \mathrm{B} _ {p, *} \rightarrow \mathrm{Z} _ {p, *} \rightarrow \mathrm{H} _ {p, *} \rightarrow 0; \quad 0 \rightarrow \mathrm{Z} _ {p, *} \rightarrow \mathrm{P} _ {p, *} \rightarrow \mathrm{B} _ {p + 1, *} \rightarrow 0).
\]

\begin{enumerate}
\def\labelenumi{\alph{enumi})}
\setcounter{enumi}{1}
\item
  设 \(\mathbf{C}'\) 为另一个复形， \((\mathbf{P}', p')\) 为 \(\mathbf{C}'\) 的一个投射分解， \(u: \mathbf{C} \to \mathbf{C}'\) 为复形的一个态射。证明存在一个双复形的态射 \(\tilde{u}: \mathbf{P} \to \mathbf{P}'\) ，使得 \(u \circ p = p' \circ \tilde{u}\) （对每个 \(p \in \mathbf{Z}\) ，选取态射 \(\mathbf{B}_{p'}', (\mathbf{P}) \to \mathbf{B}_{p'}', (\mathbf{P}')\) 与 \(\mathrm{H}_{p'}', (\mathbf{P}) \to \mathrm{H}_{p'}', (\mathbf{P}')\) ；如前所述构造 \(\tilde{u}\) ，使用习题 16，b)）。
\item
  设 \(v: \mathbf{C} \to \mathbf{C}'\) 为与 \(u\) 同伦的复形的态射，并设 \(\tilde{v}: \mathbf{P} \to \mathbf{P}'\) 为一个双复形的态射，使得 \(v \circ p = p' \circ \tilde{v}\) 。证明 \(\tilde{u}\) 与 \(\tilde{v}\) 是同伦的（按 X，第 174 页，习题 11 c) 的意义；利用习题 16 c)）。
\end{enumerate}

\(d)\) 设 E 为 \(\mathbf{Z}\) 的一个子集，且 C 为一个 A-复形，使得 \(\mathbf{C}_{p}=0\) （当 \(p\in\mathbf{E}\) ）。证明存在 C 的一个投射分解 \((\mathbf{P},p)\) ，使得 \(\mathbf{P}_{p,q}=0\) （当 \(p\in\mathbf{E}\) ）。

\begin{enumerate}
\def\labelenumi{\alph{enumi})}
\setcounter{enumi}{4}
\item
  假设存在一个整数 \(n\) ，使得每个 A-模都容许一个长度为 \(\leqslant n\) 的投射分解。证明每个 A-复形 C 都容许一个投射分解 \((P, p)\) ，使得 \(P_{p,q} = 0\) （当 \(q > n\) ）。
\item
  设 C 是一个 A-复形。C 的一个内射分解是一个二元组 \((e, l)\) ，其中 I 是一个双复形且 \(e: C \to 1\) 是双复形的一个态射，使得 \(I^{p,q} = 0\) （对 q \textless{} 0）且使得复形 \(I^{p,*}\) （相应地， \(Z^{p,*}(I)\) , \(B^{p,*}(l)\) , \(H^{p,*}(l)\) ）定义了 \(C^{p}\) （相应地， \(Z^{p}(C)\) , \(B^{p}(C)\) , \(H^{p}(C)\) ）的一个内射分解（对每个 \(p \in Z\) ）。陈述并证明内射分解情形下结果 a) 到 e) 的类比。
\end{enumerate}

\begin{enumerate}
\def\labelenumi{\arabic{enumi})}
\setcounter{enumi}{17}
\tightlist
\item
  假设环 A 是交换的。一个微分分次 A-代数是指一个分次类型为 Z 的 A-代数 S，使得 \(S_{n} = 0\) （对 n \textless{} 0），并配备一个次数为 (-1) 的反导子 d，满足 \(d \circ d = 0\) 。我们也用 S 表示 A-模复形 (S, d)。
\end{enumerate}

\begin{enumerate}
\def\labelenumi{\alph{enumi})}
\item
  设 C 为 A-模复形，在次数 \textless{} 0 处为零。证明在 S = T(C)（相应地 S = S(C)，相应地 S = Λ(C)）上存在一个微分分次 A-代数结构，使得 C 到 S 的典范嵌入是复形的态射。
\item
  证明 S 的乘法在 H(S) 上诱导一个分次 A-代数结构。
\item
  设 S、T 为两个微分分次 A-代数。令 D(s ⊗ t) = ds ⊗ t + (-1)\^{}p s ⊗ dt（对 s ∈ S\_p, t ∈ T）。证明 D 在左张量积 S *⊗ \_A T（III，第 49 页）上定义了微分分次 A-代数的结构。
\end{enumerate}

\begin{enumerate}
\def\labelenumi{\arabic{enumi})}
\setcounter{enumi}{18}
\tightlist
\item
  假设环 A 是交换的。
\end{enumerate}

\begin{enumerate}
\def\labelenumi{\alph{enumi})}
\item
  设 \((\mathbb{S}, d_{\mathbb{S}})\) 为一个微分分次 A-代数（习题 18），M 是一个 A-模， \(u: \mathbb{M} \to Z_{n}(\mathbb{S})\) 为一个 A-同态。证明在 \(\mathbb{S} \otimes_{\mathbb{A}} \mathbf{T}(\mathbb{M})\) 上存在一个微分分次 A-代数结构，使得 \(d(s \otimes 1) = d_{\mathbb{S}} s \otimes 1\) （当 \(s \in \mathbb{S}\) ）且 \(d(1 \otimes m) = u(m) \otimes 1\) （当 \(m \in \mathbb{M}\) ）。
\item
  设 B 是一个 A-代数。证明存在一个微分分次 A-代数 S 和一个 A-代数同态 \(p: \mathbf{S}_{0} \to \mathbf{B}\) ，使得 \((\mathbf{S}, p)\) 是 A-模 B 的一个自由分解。（通过对 \(n\) 作归纳，利用 \(a\) )，构造一个微分分次代数 \(\mathbf{S}(n)\) ，它在每个次数都是自由的，使得 \(\mathrm{H}_{i}(\mathbf{S}(n)) = 0\) （当 \(0 < i < n\) ）且 \(\mathrm{H}_{0}(\mathbf{S}(n)) = \mathbf{B}\) 。）
\item
  若代数 B 是交换的，证明可以找到一个交错 A-代数 S（III，第 53 页）满足 b) 的条件。若进一步 A 是诺特的，且 B 是 A 关于一个理想的商，证明可以找到 S 使得 \(S_{0} = A\) 并且 \(S_{n}\) 对每个 n 都是有限生成的 A-模。
\end{enumerate}

¶ 20) 假设环 A 是交换环 k 上的一个代数。设 \(\eta: k \to A\) 为由 \(\eta(\lambda) = \lambda 1_{A}\) （当 \(\lambda \in k\) ）定义的同态，并设 \(\overline{A}\) 为其余核。若 \(a\) 是 A 的一个元素，我们用 \(\overline{a}\) 表示它在 \(\overline{A}\) 中的类。我们设 \(\mathrm{N}_{n}(\mathrm{A}) = \mathrm{A} \otimes_{k} \overline{\mathrm{A}}^{\otimes n} \otimes_{k} \mathrm{A}\) （当 \(n \geqslant 0\) ）, \(\mathrm{N}_{n}(\mathrm{A}) = 0\) （当 \(n < 0\) ），且

\[
\mathrm{N} (\mathbf {A}) = \bigoplus_ {n \in \mathbf {Z}} \mathrm{N} _ {n} (\mathbf {A}).
\]

\begin{enumerate}
\def\labelenumi{\alph{enumi})}
\tightlist
\item
  证明存在 N(A) 的一个分次次数为 (-1) 的 (A, A)-自同态 d，使得
\end{enumerate}

\[
\begin{array}{l} d (a \otimes \bar {a} _ {1} \otimes \dots \otimes \bar {a} _ {n} \otimes b) = a a _ {1} \otimes \bar {a} _ {2} \otimes \dots \otimes \bar {a} _ {n} \otimes b + \sum_ {i = 1} ^ {n - 1} (- 1) ^ {i} a \otimes \dots \otimes \overline {{a _ {i}}} \overline {{a _ {i + 1}}} \otimes \dots \otimes \bar {a} _ {n} \otimes b \\ \qquad + (- 1) ^ {n} a \otimes \bar {a} _ {1} \otimes \dots \otimes \bar {a} _ {n - 1} \otimes a _ {n} b \quad \text {for} \quad a, a _ {1},..., a _ {n}, b \in A \end{array}
\]

\begin{enumerate}
\def\labelenumi{\alph{enumi})}
\setcounter{enumi}{1}
\item
  证明 \(d \circ d = 0\) ，使得 \(\mathbf{N}(\mathbf{A})\) 是 \((\mathbf{A}, \mathbf{A})\) -双模的一个复形。若 M 是一个左 A-模，我们记 \(\mathbf{N}(\mathbf{A}, \mathbf{M})\) 为 A-复形 \(\mathbf{N}(\mathbf{A}) \otimes_{\mathbf{A}} \mathbf{M}\) 。
\item
  设 B(A, M) 为 M 的标准分解（X，第 58 页）， \(p: \mathrm{B}(\mathrm{A}, \mathrm{M}) \to \mathrm{N}(\mathrm{A}, \mathrm{M})\) 为由 \(p(a_0 \otimes \cdots \otimes a_n \otimes m) = a_0 \otimes \overline{a}_1 \otimes \cdots \otimes \overline{a}_n \otimes m\) （当 \(a_0, \ldots, a_n \in \mathbf{A}, m \in \mathbf{M}\) ）定义的 A-同态。证明 \(p\) 是 A-复形的一个同伦等价（通过归纳构造一个在同伦意义下的逆映射）。
\end{enumerate}

* 21) 设 \(k\) 是一个交换环，并设 g 是一个 \(k\) 李代数，U 是其包络代数（参见 LIE, I, § 2）。假设 g 是一个 \(k\) -自由模。

\begin{enumerate}
\def\labelenumi{\alph{enumi})}
\tightlist
\item
  证明对每个 \(n \geqslant 0\) 存在一个单射 U-同态
\end{enumerate}

\[
j _ {n}: \mathrm{U} \otimes_ {k} \Lambda^ {n} (\mathrm{g}) \rightarrow \mathrm{U} ^ {\otimes (n + 1)} \quad \text { such that } \quad j _ {n} (u \otimes (x _ {1} \wedge \dots \wedge x _ {n})) = \sum_ {\sigma \in \mathfrak {S} _ {n}} \varepsilon (\sigma), u \otimes x _ {\sigma (1)} \otimes \dots \otimes x _ {\sigma (n)}
\]

当 \(u\in \mathbf{U},x_1,\dots ,x_n\in \mathfrak{g}\)

因此我们定义一个次数为 0 的分次 U-同态 \(j:U\otimes_{k}\Lambda(g)\to B(U,k)\) ，其中 \(B(U,k)\) 是左 U-模 \(k\) 的标准分解（ \(k\) 的 U-模结构由自然同态 \(U\to k\) 导出）。证明 \(U\otimes_{k}\Lambda(g)\) 通过 \(j\) 等同于 \(B(U,k)\) 的一个子复形，其微分由下述公式给出：

\[
\begin{array}{l} d (u \otimes (x _ {1} \wedge \dots \wedge x _ {n})) = \sum_ {1 \leqslant i \leqslant n} (- 1) ^ {i + 1} u x _ {i} \otimes (x _ {1} \wedge \dots \wedge \hat {x} _ {i} \wedge \dots \wedge x _ {n}) \\ + \sum_ {1 \leqslant i <   j \leqslant n} (- 1) ^ {i + j} u \otimes ([ x _ {i}, x _ {j} ] \wedge x _ {1} \wedge \dots \wedge \hat {x} _ {i} \wedge \dots \wedge \hat {x} _ {j} \wedge \dots \wedge x _ {n}) \quad \text { pour } \quad u \in U, x _ {1},..., x _ {n} \in g. \end{array}
\]

（字母上方的符号表示该字母应被省略）。我们用 V(g) 表示如此定义的 U-复形。b) 证明 V(g) 定义了 U-模 k 的一个自由分解（令 \((\mathrm{U}_{n})_{n\geqslant0}\) 为 U 的自然滤过， \(F_{r}V(g)\) 为 V(g) 的由元素 \(u_{p}\otimes x_{q}\) （当 \(u_{p}\in U_{p}, x_{q}\in\Lambda^{q}(g), p+q\leqslant r\) ）生成的子 k-模。证明 \(F_{r}V(g)\) 是 V(g) 的一个子复形，并且复形 \(\bigoplus_{r\geqslant0}F_{r}V(g)/F_{r-1}V(g)\)

同构于复形 \(\mathbf{S}(\mathfrak{g})\otimes\symbf{\Lambda}(\mathfrak{g})\) （在 X，第 151 页中定义）。*
